\chapter{Geometric Measure Theory}\label{geometric-measure-theory}
\chapterstatus{complete}

\section{Introduction}\label{geometric-measure-theory:section-introduction}

Lebesgue measure measures $n$-dimensional volume in $\RR^n$; geometric measure theory measures sets of lower, possibly fractional, dimension: curves and surfaces, but also the Cantor set, the Koch
curve and the Sierpi\'nski gasket. This chapter begins with the Brunn--Minkowski inequality for the volume of sums of sets, with its consequences the isodiametric and the isoperimetric
inequalities. It then introduces Hausdorff measures and Hausdorff dimension, with the Cantor set as first example, the box-counting dimension, and self-similar sets generated by
iterated function systems, whose dimension is computed by Moran's equation (Theorem~\ref{geometric-measure-theory:theorem-moran}). It then proves Rademacher's theorem on the differentiability of Lipschitz
maps, the area formula, which computes the $k$-dimensional measure of parametrised sets, the coarea formula with integration in polar coordinates, and the first results on densities and rectifiable sets.
References: \cite{falconer-fractal}, \cite{mattila-1995}, \cite{evans-gariepy} and \cite{federer-gmt}.

\noindent\textbf{Prerequisites.} Chapter~\ref{general-topology} (General Topology), Chapter~\ref{real-analysis} (Real Analysis), Chapter~\ref{measure-theory} (Measure and Integration),
Chapter~\ref{lebesgue-differentiation} (Differentiation and Absolute Continuity) and Chapter~\ref{radon-measures} (Radon Measures).

\section{The Brunn--Minkowski inequality}\label{geometric-measure-theory:section-brunn-minkowski}

Among all plane regions of given perimeter, the disc has the largest area: this is the \emph{isoperimetric problem}, known to the Greeks and proved rigorously only in the nineteenth century. Its
higher-dimensional version, and the related \emph{isodiametric inequality} (among sets of given diameter, the ball has the largest volume), follow from one inequality about the volume of sums of sets,
due to Brunn and Minkowski. We prove it first, since the isodiametric inequality is needed to identify Hausdorff measure with Lebesgue measure in the next section.

\begin{notation}\label{geometric-measure-theory:notation-minkowski-sum}
For $A, B \subseteq \RR^n$ and $t \in \RR$, the \emph{Minkowski sum} is $A + B = \{a + b : a \in A, b \in B\}$, and $tA = \{ta : a \in A\}$. We write $\lambda = \lambda^n$, $\overline{B}_r = \overline{B}(0, r)$
for the closed ball, $\omega_n = \lambda(\overline{B}_1)$ for the volume of the unit ball, and $\alpha(n) = 2^{-n}\omega_n$ for the volume of a ball of diameter $1$ (the
normalising constant of Hausdorff measure in the next section). A \emph{box} is a set $\prod_i[a_i, b_i]$ with $a_i < b_i$.
\end{notation}

\begin{example}\label{geometric-measure-theory:example-minkowski-sum}
\begin{enumerate}
\item $A + \overline{B}_r$ is the closed $r$-neighbourhood $\{x : \operatorname{dist}(x, A) \leq r\}$ of a closed set $A$.
\item The sum of two boxes is a box: $\prod_i[a_i, b_i] + \prod_i[c_i, d_i] = \prod_i[a_i + c_i, b_i + d_i]$. The sum of a horizontal and a vertical unit segment in $\RR^2$ is the unit square: two null sets
can have a sum of positive measure.
\item For convex $A$ and $s, t \geq 0$, $sA + tA = (s + t)A$; for $A = \{0, 1\}$, $A + A = \{0, 1, 2\} \neq 2A$.
\end{enumerate}
\end{example}

The naive guess $\lambda(A + B) \geq \lambda(A) + \lambda(B)$ is true but weak. Since $\lambda(tA) = t^n\lambda(A)$, the natural quantity is $\lambda(A)^{1/n}$, the side of a cube with the same volume as $A$; for
homothetic convex sets $A$ and $tA$ it is additive: $\lambda(A + tA)^{1/n} = (1 + t)\lambda(A)^{1/n}$.

\begin{theorem}[Brunn--Minkowski inequality]\label{geometric-measure-theory:theorem-brunn-minkowski}
For nonempty compact sets $A, B \subseteq \RR^n$,
\begin{equation}\label{geometric-measure-theory:equation-brunn-minkowski}
\lambda(A + B)^{1/n} \geq \lambda(A)^{1/n} + \lambda(B)^{1/n}.
\end{equation}
The same holds for Lebesgue measurable $A$, $B$ such that $A + B$ is Lebesgue measurable.
\end{theorem}

\begin{proof}
$A + B$ is compact, as the image of $A \times B$ under addition.

\emph{Step 1: two boxes.} Let $A$ and $B$ be boxes with sides $a_i$ and $b_i$; then $A + B$ is a box with sides $a_i + b_i$. Dividing \eqref{geometric-measure-theory:equation-brunn-minkowski} by
$\prod_i(a_i + b_i)^{1/n}$, it becomes
\[
\prod_i\Bigl(\frac{a_i}{a_i + b_i}\Bigr)^{1/n} + \prod_i\Bigl(\frac{b_i}{a_i + b_i}\Bigr)^{1/n} \leq \frac1n\sum_i\frac{a_i}{a_i + b_i} + \frac1n\sum_i\frac{b_i}{a_i + b_i} = 1,
\]
where we used the inequality between the geometric and the arithmetic mean (Example~\ref{measure-theory:example-jensen}(3)) twice.

\emph{Step 2: finite unions of boxes.} Let $A$ and $B$ be finite unions of boxes with pairwise disjoint interiors, $k$ boxes in all; we induct on $k \geq 2$, the case $k = 2$ being Step 1. Suppose that
$A$ consists of at least two boxes (otherwise exchange $A$ and $B$). Two boxes with disjoint interiors have disjoint interiors of some coordinate projection, so there is a hyperplane
$H = \{x_i = c\}$ with one box of $A$ in $\{x_i \leq c\}$ and another in $\{x_i \geq c\}$. Translating, which changes neither side of the inequality, we may assume $c = 0$. Let
$A^+ = A \cap \{x_i \geq 0\}$ and $A^- = A \cap \{x_i \leq 0\}$, where we discard the pieces of zero volume (faces lying in $H$), which changes no measure and only shrinks the sums below; each is a union
of boxes with disjoint interiors, with fewer boxes than $A$, since each box of $A$ contributes at most one box to $A^\pm$ and one box contributes none. Translating $B$ in the $i$-th coordinate direction, the function $s \mapsto \lambda((B + se_i) \cap \{x_i \geq 0\})$ is continuous and goes from $0$ to $\lambda(B)$; choose the
translate for which $B^\pm = B \cap \{\pm x_i \geq 0\}$ satisfy
\[
\frac{\lambda(B^+)}{\lambda(B)} = \frac{\lambda(A^+)}{\lambda(A)} =: \theta \in (0, 1), \qquad \frac{\lambda(B^-)}{\lambda(B)} = 1 - \theta.
\]
Each $B^\pm$ has at most as many boxes as $B$, so the pairs $(A^+, B^+)$ and $(A^-, B^-)$ have fewer than $k$ boxes in all. The sets $A^+ + B^+ \subseteq \{x_i \geq 0\}$ and $A^- + B^- \subseteq \{x_i \leq 0\}$
are contained in $A + B$ and overlap in a null set, so by the induction hypothesis
\[\begin{aligned}
\lambda(A + B) &\geq \lambda(A^+ + B^+) + \lambda(A^- + B^-) \\
&\geq \bigl(\lambda(A^+)^{1/n} + \lambda(B^+)^{1/n}\bigr)^n + \bigl(\lambda(A^-)^{1/n} + \lambda(B^-)^{1/n}\bigr)^n \\
&= \theta\bigl(\lambda(A)^{1/n} + \lambda(B)^{1/n}\bigr)^n + (1 - \theta)\bigl(\lambda(A)^{1/n} + \lambda(B)^{1/n}\bigr)^n \\
&= \bigl(\lambda(A)^{1/n} + \lambda(B)^{1/n}\bigr)^n.
\end{aligned}\]

\emph{Step 3: compact sets.} For $j \geq 1$ let $A_j$ be the union of the closed dyadic cubes of side $2^{-j}$ that meet $A$, and similarly $B_j$; these are finite unions of boxes with disjoint
interiors, $A_j \supseteq A$, and $A_j \subseteq A + \overline{B}_{\varepsilon_j}$ with $\varepsilon_j = \sqrt n\,2^{-j}$. The compact sets $A_j$ decrease, as each cube of side $2^{-j-1}$ meeting $A$ lies in a
cube of side $2^{-j}$ meeting $A$, and their intersection is $A$, since $A$ is closed; so $\lambda(A_j) \to \lambda(A)$ by continuity from above
(Proposition~\ref{measure-theory:proposition-measure-properties}), and likewise for $B$. Moreover $A_j + B_j \subseteq (A + B) + \overline{B}_{2\varepsilon_j}$, and these neighbourhoods decrease to the compact set
$A + B$. By Step 2,
\[\begin{aligned}
\lambda(A + B) &= \lim_j\lambda\bigl((A + B) + \overline{B}_{2\varepsilon_j}\bigr) \geq \limsup_j\bigl(\lambda(A_j)^{1/n} + \lambda(B_j)^{1/n}\bigr)^n \\
&= \bigl(\lambda(A)^{1/n} + \lambda(B)^{1/n}\bigr)^n.
\end{aligned}\]

\emph{Step 4: measurable sets.} If $A$, $B$ and $A + B$ are measurable and nonempty, compact sets $K \subseteq A$ and $L \subseteq B$ satisfy $K + L \subseteq A + B$, and by inner regularity
(Theorem~\ref{measure-theory:theorem-lebesgue-measure}) $\lambda(K)$ and $\lambda(L)$ can be taken arbitrarily close to $\lambda(A)$ and $\lambda(B)$, or arbitrarily large if these are infinite.
\end{proof}

\begin{remark}\label{geometric-measure-theory:remark-brunn-minkowski}
\begin{enumerate}
\item Sums can be much larger than their summands: for the Cantor set $C$, $C + C = [0, 2]$. Indeed $\frac12C$ consists of the numbers $\sum_i\varepsilon_i3^{-i}$ with $\varepsilon_i \in \{0, 1\}$,
and the digitwise sums $\varepsilon_i + \varepsilon_i^{\prime}$ run through all ternary digits $\{0, 1, 2\}$, so $\frac12C + \frac12C = [0, 1]$. Here the inequality reads $2 \geq 0 + 0$. For measurable sets
$A$, $B$ nothing in our arguments shows that $A + B$ is measurable, which is why its measurability is assumed in the last statement.
\item Equality holds in \eqref{geometric-measure-theory:equation-brunn-minkowski} for compact sets of positive measure if and only if $A$ and $B$ are, up to null sets, homothetic convex bodies \cite{gardner-2002}.
For instance, for two segments of different directions in $\RR^2$ the left side is positive and the right side is $0$.
\item The exponent $1/n$ is the best possible for an inequality of the form $\lambda(A + B)^p \geq \lambda(A)^p + \lambda(B)^p$: for $B = tA$ with $A$ convex, $\lambda(A + B) = (1 + t)^n\lambda(A)$, and
$(1 + t)^{np} \geq 1 + t^{np}$ fails for small $t$ when $np > 1$. For $p < 1/n$ the inequality follows from the case $p = 1/n$, by concavity of $s \mapsto s^{pn}$.
\end{enumerate}
\end{remark}

\begin{corollary}\label{geometric-measure-theory:corollary-brunn-minkowski-multiplicative}
For nonempty compact $A, B \subseteq \RR^n$ and $0 < t < 1$, $\lambda((1 - t)A + tB) \geq \lambda(A)^{1 - t}\lambda(B)^t$. In particular $t \mapsto \lambda((1 - t)A + tB)^{1/n}$ is concave when $A$ and $B$ are
convex.
\end{corollary}

\begin{proof}
By the theorem and the homogeneity $\lambda(sA) = s^n\lambda(A)$, $\lambda((1 - t)A + tB)^{1/n} \geq (1 - t)\lambda(A)^{1/n} + t\lambda(B)^{1/n}$, and the weighted inequality between arithmetic and geometric means
(concavity of $\log$) gives $(1 - t)a + tb \geq a^{1 - t}b^t$ for $a, b \geq 0$. For convex $A$, $B$ and $s, t, u \in [0, 1]$ with $u = (1 - \theta)s + \theta t$,
$(1 - u)A + uB = (1 - \theta)\bigl((1 - s)A + sB\bigr) + \theta\bigl((1 - t)A + tB\bigr)$, and the first inequality of this proof gives concavity.
\end{proof}

\begin{theorem}[Isodiametric inequality]\label{geometric-measure-theory:theorem-isodiametric}
For every bounded $C \subseteq \RR^n$, $\lambda^*(C) \leq \omega_n\bigl(\tfrac12\operatorname{diam}C\bigr)^n = \alpha(n)(\operatorname{diam}C)^n$: among all sets of given diameter, the ball has the
largest volume.
\end{theorem}

\begin{proof}
The closure $A = \overline{C}$ is compact with the same diameter $d$ and $\lambda^*(C) \leq \lambda(A)$, so we may assume $C = A$ compact and nonempty. The set $-A$ has the same measure as $A$, and
$\frac12A + \frac12(-A) = \{\frac12(x - y) : x, y \in A\}$ lies in the ball $\overline{B}_{d/2}$, since $\|x - y\| \leq d$. By Brunn--Minkowski,
\[
\omega_n(d/2)^n \geq \lambda\bigl(\tfrac12A + \tfrac12(-A)\bigr) \geq \bigl(\tfrac12\lambda(A)^{1/n} + \tfrac12\lambda(-A)^{1/n}\bigr)^n = \lambda(A).
\]
\end{proof}

The inequality is not obvious, because a set of diameter $d$ need not lie in a ball of diameter $d$: an equilateral triangle of side $1$ has diameter $1$, but the smallest disc containing it has
radius $1/\sqrt3 > \frac12$. Yet its area $\sqrt3/4 \approx 0.433$ is less than $\pi/4 \approx 0.785$.

\subsection{The isoperimetric inequality}\label{geometric-measure-theory:subsection-isoperimetric}

Measuring the boundary of an arbitrary set is delicate. The simplest notion uses the volume of thin neighbourhoods: the shell $(A + \overline{B}_\varepsilon) \setminus A$ has volume about
$\varepsilon$ times the area of the boundary.

\begin{definition}\label{geometric-measure-theory:definition-minkowski-content}
The \emph{outer Minkowski content} of a compact set $A \subseteq \RR^n$ is
\[
\mathcal{M}(A) = \liminf_{\varepsilon \to 0^+}\frac{\lambda(A + \overline{B}_\varepsilon) - \lambda(A)}{\varepsilon}.
\]
\end{definition}

\begin{example}\label{geometric-measure-theory:example-minkowski-content}
\begin{enumerate}
\item For the ball $\overline{B}_r$, $\lambda(\overline{B}_r + \overline{B}_\varepsilon) = \omega_n(r + \varepsilon)^n$, and $\mathcal{M}(\overline{B}_r) = n\omega_nr^{n-1}$, the area of the sphere of radius
$r$ (for $n = 2$, the circumference $2\pi r$; for $n = 3$, $4\pi r^2$).
\item \emph{Steiner's formula for convex polygons.} Let $P \subseteq \RR^2$ be a convex polygon with perimeter $L$. Then $P + \overline{B}_\varepsilon$ is the union of $P$, of a rectangle of height $\varepsilon$
on each side (of total area $L\varepsilon$), and of a circular sector of radius $\varepsilon$ at each vertex, with angle $\pi$ minus the interior angle; the exterior angles add up to $2\pi$, so the
sectors form a full disc, and these pieces overlap only on segments. Hence $\lambda(P + \overline{B}_\varepsilon) = \lambda(P) + L\varepsilon + \pi\varepsilon^2$ and $\mathcal{M}(P) = L$.
\item For a segment $S$ of length $\ell$ in $\RR^2$, $\lambda(S + \overline{B}_\varepsilon) = 2\ell\varepsilon + \pi\varepsilon^2$, so $\mathcal{M}(S) = 2\ell$: the outer Minkowski content sees both sides of a set
without interior.
\end{enumerate}
\end{example}

\begin{theorem}[Isoperimetric inequality]\label{geometric-measure-theory:theorem-isoperimetric}
For every compact $A \subseteq \RR^n$,
\[
\mathcal{M}(A) \geq n\omega_n^{1/n}\lambda(A)^{(n-1)/n},
\]
with equality for balls. Equivalently: a compact set has outer Minkowski content at least that of the ball of the same volume.
\end{theorem}

\begin{proof}
By Brunn--Minkowski and the binomial theorem, keeping only the first two terms,
\[
\lambda(A + \overline{B}_\varepsilon) \geq \bigl(\lambda(A)^{1/n} + \varepsilon\omega_n^{1/n}\bigr)^n \geq \lambda(A) + n\varepsilon\omega_n^{1/n}\lambda(A)^{(n-1)/n}.
\]
Subtract $\lambda(A)$, divide by $\varepsilon$ and let $\varepsilon \to 0$. For the ball $\overline{B}_r$, $n\omega_n^{1/n}(\omega_nr^n)^{(n-1)/n} = n\omega_nr^{n-1} = \mathcal{M}(\overline{B}_r)$ by
Example~\ref{geometric-measure-theory:example-minkowski-content}(1). If $\lambda(A) = \lambda(\overline{B}_r)$, the right side is the content of $\overline{B}_r$.
\end{proof}

\begin{corollary}\label{geometric-measure-theory:corollary-isoperimetric-plane}
A convex polygon of perimeter $L$ and area $F$ satisfies $L^2 \geq 4\pi F$.
\end{corollary}

\begin{proof}
For $n = 2$ the theorem reads $\mathcal{M}(P) \geq 2\sqrt\pi F^{1/2}$, and $\mathcal{M}(P) = L$ by Example~\ref{geometric-measure-theory:example-minkowski-content}(2).
\end{proof}

For a square of side $1$, $L^2 = 16 > 4\pi$; for the regular $m$-gon inscribed in the unit circle, $L = 2m\sin(\pi/m)$ and $F = \frac m2\sin(2\pi/m)$, and $L^2/(4\pi F) \to 1$ as $m \to \infty$: the
inequality is sharp, but no polygon attains it. For a compact $(n-1)$-rectifiable set $S$, such as the boundary of a set with smooth boundary,
$\lambda(S + \overline{B}_\varepsilon)/(2\varepsilon) \to \mathcal{H}^{n-1}(S)$ \cite[3.2.39]{federer-gmt}; for a set $A$ with smooth boundary the outer half of this shell is $(A + \overline{B}_\varepsilon) \setminus A$,
so $\mathcal{M}(A) = \mathcal{H}^{n-1}(\partial A)$, and the theorem becomes the classical $\mathcal{H}^{n-1}(\partial A) \geq n\omega_n^{1/n}\lambda(A)^{(n-1)/n}$; its relation with the Sobolev inequality is described in Remark~\ref{geometric-measure-theory:remark-coarea}.
A different proof in the plane, by Fourier series, is Hurwitz's (Chapter~\ref{elliptic-operators}).

\section{Hausdorff measures}\label{geometric-measure-theory:section-hausdorff}

Lebesgue measure $\lambda^n$ measures $n$-dimensional volume in $\RR^n$, but it does not see lower-dimensional sets: a curve in the plane, or a surface in space, has measure $0$. How long is
a curve, and what is the area of a surface, without choosing a parametrisation? And how big is the Cantor set, which has measure $0$ and is uncountable? Hausdorff measures answer these
questions at once: for every $s \geq 0$ there is an $s$-dimensional measure $\mathcal{H}^s$ on any metric space, defined by covering a set with small sets $C$ and adding up
$(\operatorname{diam}C)^s$. This and the following sections are an introduction to \emph{geometric measure theory}: Hausdorff measure and dimension, self-similar sets, differentiability of
Lipschitz maps, the area formula, which computes $\mathcal{H}^k$ of parametrised $k$-dimensional sets, the coarea formula, which slices sets by level sets, and densities and rectifiable sets. We follow \cite[Chapters 2--3]{evans-gariepy}; the comprehensive reference is
\cite{federer-gmt}.

Throughout, $X$ is a metric space, $\operatorname{diam}C = \sup\{d(x, y) : x, y \in C\}$ for $C \neq \emptyset$, and $\operatorname{diam}\emptyset = 0$. For $s \geq 0$ let
$\alpha(s) = \pi^{s/2}/(2^s\Gamma(\frac s2 + 1))$, with $\Gamma$ as in Exercise~\ref{measure-theory:exercise-ball-volume}; for an integer $n$, $\alpha(n)$ is the volume of a ball of
diameter $1$ in $\RR^n$. The constant only serves to make $\mathcal{H}^n$ equal to Lebesgue measure on $\RR^n$ (Theorem~\ref{geometric-measure-theory:theorem-hausdorff-lebesgue}); in the sums below
the empty set contributes $0$ and every nonempty set $C$ contributes $\alpha(s)(\operatorname{diam}C)^s$, with the convention $(\operatorname{diam}C)^0 = 1$, also when $C$ is a
single point.

\begin{definition}\label{geometric-measure-theory:definition-hausdorff}
For $A \subseteq X$, $s \geq 0$ and $\delta > 0$ let
\[
\mathcal{H}^s_\delta(A) = \inf\Bigl\{\sum_k\alpha(s)(\operatorname{diam}C_k)^s : A \subseteq \bigcup_kC_k,\ \operatorname{diam}C_k \leq \delta\Bigr\},
\]
the infimum over countable covers by sets of diameter at most $\delta$. As $\delta$ decreases, fewer covers are allowed and $\mathcal{H}^s_\delta(A)$ increases; the
\emph{$s$-dimensional Hausdorff (outer) measure} is $\mathcal{H}^s(A) = \lim_{\delta \to 0}\mathcal{H}^s_\delta(A) = \sup_{\delta > 0}\mathcal{H}^s_\delta(A)$.
\end{definition}

The restriction $\operatorname{diam}C_k \leq \delta$ is essential: with arbitrary covers, a long segment could be covered by a single set and would get the same size as a short one, while
small sets are forced to follow the shape of $A$. Replacing each $C_k$ by its closure does not change its diameter, so closed covers give the same infimum. And $\mathcal{H}^s(A)$ only
depends on the metric space $A$: a cover of $A$ by subsets of $X$ gives the cover by the sets $C_k \cap A$, of smaller diameters. In particular an isometry $A \to B$ between subsets of two
metric spaces preserves $\mathcal{H}^s$.

\begin{lemma}[Carath\'eodory's criterion]\label{geometric-measure-theory:lemma-metric-outer-measure}
Let $\mu^*$ be an outer measure on $X$ which is \emph{metric}: $\mu^*(A \cup B) = \mu^*(A) + \mu^*(B)$ whenever $d(A, B) = \inf\{d(a, b) : a \in A, b \in B\} > 0$. Then every Borel set is
$\mu^*$-measurable.
\end{lemma}

\begin{proof}
The measurable sets form a $\sigma$-algebra (Theorem~\ref{measure-theory:theorem-caratheodory}), so it suffices to show that a closed set $F$ is measurable, that is,
$\mu^*(E \cap F) + \mu^*(E \setminus F) \leq \mu^*(E)$ for every $E$ with $\mu^*(E) < \infty$. Let $E_j = \{x \in E : d(x, F) \geq 1/j\}$. Since $d(E \cap F, E_j) \geq 1/j$,
$\mu^*(E \cap F) + \mu^*(E_j) = \mu^*((E \cap F) \cup E_j) \leq \mu^*(E)$, and it remains to show $\mu^*(E_j) \to \mu^*(E \setminus F)$. Let
$R_j = \{x \in E : \frac1{j+1} \leq d(x, F) < \frac1j\}$. As $F$ is closed, $E \setminus F = E_j \cup \bigcup_{i \geq j}R_i$. If $i \geq l + 2$, then $d(R_i, R_l) > 0$, so by induction
$\sum_{i \leq N}\mu^*(R_{2i}) = \mu^*(\bigcup_{i \leq N}R_{2i}) \leq \mu^*(E)$, and similarly for the odd indices. Hence $\sum_i\mu^*(R_i) < \infty$, and
$\mu^*(E \setminus F) \leq \mu^*(E_j) + \sum_{i \geq j}\mu^*(R_i)$, where the tail tends to $0$. With $\mu^*(E_j) \leq \mu^*(E \setminus F)$ this gives the limit.
\end{proof}

This is \emph{Carath\'eodory's criterion}; it is a different statement from Carath\'eodory's theorem (Theorem~\ref{measure-theory:theorem-caratheodory}), which it uses: the theorem builds a measure
from any outer measure, and the criterion says that for a metric outer measure the resulting $\sigma$-algebra contains the Borel sets.

\begin{proposition}\label{geometric-measure-theory:proposition-hausdorff}
\begin{enumerate}
\item $\mathcal{H}^s$ is a metric outer measure, so its restriction to the Borel sets is a measure.
\item $\mathcal{H}^0$ is the counting measure.
\item If $f \colon A \to Y$ is Lipschitz with constant $L$, then $\mathcal{H}^s(f(A)) \leq L^s\mathcal{H}^s(A)$. For $X = \RR^n$, $\mathcal{H}^s$ is translation invariant and
$\mathcal{H}^s(rA) = r^s\mathcal{H}^s(A)$ for $r > 0$.
\item If $\mathcal{H}^s(A) < \infty$ and $t > s$, then $\mathcal{H}^t(A) = 0$. If $\mathcal{H}^t(A) > 0$ and $s < t$, then $\mathcal{H}^s(A) = \infty$.
\end{enumerate}
\end{proposition}

\begin{proof}
(1) Each $\mathcal{H}^s_\delta$ is an outer measure: it is monotone, and given covers of the $A_k$ with sums at most $\mathcal{H}^s_\delta(A_k) + \varepsilon2^{-k}$, their union covers
$\bigcup_kA_k$. Then $\mathcal{H}^s(\bigcup_kA_k) = \sup_\delta\mathcal{H}^s_\delta(\bigcup_kA_k) \leq \sup_\delta\sum_k\mathcal{H}^s_\delta(A_k) \leq \sum_k\mathcal{H}^s(A_k)$. If
$d(A, B) > 0$ and $\delta < d(A, B)$, no set of diameter at most $\delta$ meets both $A$ and $B$, so a $\delta$-cover of $A \cup B$ splits into covers of $A$ and of $B$, and
$\mathcal{H}^s_\delta(A \cup B) \geq \mathcal{H}^s_\delta(A) + \mathcal{H}^s_\delta(B)$; let $\delta \to 0$. Apply Lemma~\ref{geometric-measure-theory:lemma-metric-outer-measure}.

(2) $\alpha(0) = 1$. A set of $N$ points needs at least $N$ nonempty sets of diameter less than the smallest distance between the points, and the $N$ singletons suffice; so
$\mathcal{H}^0_\delta$ of it is $N$ for small $\delta$. An infinite set contains finite sets of every size.

(3) If $L = 0$, then $f$ is constant, $f(A)$ has at most one point, and the inequality is clear (with $0^0 = 1$ when $s = 0$). Let $L > 0$. If $(C_k)$ is a
$\delta$-cover of $A$, then $(f(C_k \cap A))$ is an $L\delta$-cover of $f(A)$ with $\operatorname{diam}f(C_k \cap A) \leq L\operatorname{diam}C_k$; so
$\mathcal{H}^s_{L\delta}(f(A)) \leq L^s\mathcal{H}^s_\delta(A)$. Translations are isometries, and $x \mapsto rx$ multiplies diameters by $r$, with inverse $x \mapsto x/r$.

(4) For a $\delta$-cover,
\[
\sum_k\alpha(t)(\operatorname{diam}C_k)^t \leq \frac{\alpha(t)}{\alpha(s)}\delta^{t-s}\sum_k\alpha(s)(\operatorname{diam}C_k)^s,
\]
so $\mathcal{H}^t_\delta(A) \leq \frac{\alpha(t)}{\alpha(s)}\delta^{t-s}\mathcal{H}^s_\delta(A) \leq \frac{\alpha(t)}{\alpha(s)}\delta^{t-s}\mathcal{H}^s(A) \to 0$. The second statement is the
contrapositive.
\end{proof}

\begin{remark}\label{geometric-measure-theory:remark-hausdorff-conventions}
\begin{enumerate}
\item \emph{Covers may not exist.} With the convention $\inf\emptyset = +\infty$, $\mathcal{H}^s_\delta(A) = \infty$ when $A$ has no countable cover by sets of diameter at most $\delta$. In a separable
metric space this never happens: if $\{x_k\}$ is a countable dense set, the balls $B(x_k, \delta/2)$ have diameter at most $\delta$ and cover $X$. In an uncountable set with the discrete metric
($d(x, y) = 1$ for $x \neq y$), sets of diameter less than $1$ are points, so an uncountable set has no countable $\delta$-cover for $\delta < 1$, and $\mathcal{H}^s$ of it is $\infty$ for every $s$,
consistently with (2) for $s = 0$. All subsets of $\RR^n$ and of manifolds are separable.
\item \emph{Borel regularity.} $\mathcal{H}^s$ is \emph{Borel regular}: every $A$ is contained in a Borel set $B$ with $\mathcal{H}^s(B) = \mathcal{H}^s(A)$. We may assume $\mathcal{H}^s(A) < \infty$. For
each $j$ choose a cover of $A$ by closed sets $C^{(j)}_k$ of diameter at most $1/j$ with $\sum_k\alpha(s)(\operatorname{diam}C^{(j)}_k)^s \leq \mathcal{H}^s_{1/j}(A) + 1/j$, and let
$B = \bigcap_j\bigcup_kC^{(j)}_k$, a Borel set containing $A$. Then $\mathcal{H}^s_{1/j}(B) \leq \mathcal{H}^s_{1/j}(A) + 1/j$ for all $j$, so $\mathcal{H}^s(B) \leq \mathcal{H}^s(A)$. On the other
hand $\mathcal{H}^s$ is in general not locally finite (Example~\ref{measure-theory:example-regular}), hence not a Radon measure on the whole space; restricted to a set of finite
$\mathcal{H}^s$-measure it is \cite[Section 1.1]{evans-gariepy}.
\end{enumerate}
\end{remark}

By (4), for each $A$ there is a unique $s_0 \in [0, \infty]$ with $\mathcal{H}^s(A) = \infty$ for $s < s_0$ and $\mathcal{H}^s(A) = 0$ for $s > s_0$; it is the \emph{Hausdorff dimension}
$\dim_{\mathcal{H}}A = \inf\{s : \mathcal{H}^s(A) = 0\}$. At $s = s_0$ the measure can be anything in $[0, \infty]$. By (3), Lipschitz maps do not increase Hausdorff dimension, and
bi-Lipschitz maps preserve it.

\begin{theorem}\label{geometric-measure-theory:theorem-hausdorff-lebesgue}
On the Borel sets of $\RR^n$, $\mathcal{H}^n = \lambda^n$. Consequently the subsets of $\RR^n$ have Hausdorff dimension at most $n$, and nonempty open subsets have dimension $n$.
\end{theorem}

\begin{proof}
We first show that $\mathcal{H}^n = c\lambda^n$ for a constant $0 < c < \infty$, and then that $c = 1$, using the isodiametric inequality and the Vitali
covering theorem.

\emph{The constant.} Let $Q = [0, 1)^n$. Covering $Q$ by $m^n$ cubes of side $1/m$ and diameter $\sqrt n/m$ gives $\mathcal{H}^n_{\sqrt n/m}(Q) \leq \alpha(n)n^{n/2}$, so
$\mathcal{H}^n(Q) < \infty$. Conversely, a set $C$ of diameter $d$ lies in a closed box with sides $d$ (each coordinate varies by at most $d$ on $C$), so $\lambda^*(C) \leq d^n$; for any cover
of $Q$, $\sum_k\alpha(n)(\operatorname{diam}C_k)^n \geq \alpha(n)\sum_k\lambda^*(C_k) \geq \alpha(n)$, so $c = \mathcal{H}^n(Q) \geq \alpha(n) > 0$.

\emph{Comparison.} $\mathcal{H}^n$ is a translation invariant measure on the Borel sets (Proposition~\ref{geometric-measure-theory:proposition-hausdorff}). For $j \geq 0$, $Q$ is the disjoint union of $2^{jn}$
translates of $[0, 2^{-j})^n$, so every half-open dyadic cube $\prod_i[a_i2^{-j}, (a_i + 1)2^{-j})$ ($a_i \in \ZZ$, $j \in \ZZ$) has $\mathcal{H}^n$-measure $c2^{-jn} = c\lambda^n$ of it; for
$j < 0$ use that such a cube is a disjoint union of $2^{-jn}$ unit cubes. Fix $m \geq 0$ and $P = [-2^m, 2^m)^n$. The dyadic cubes contained in $P$, together with $P$, form a family closed under
finite intersections, on which $\mathcal{H}^n$ and $c\lambda^n$ agree and are finite. By Lemma~\ref{measure-theory:lemma-dynkin} the two measures agree on the $\sigma$-algebra it generates, which
contains every open subset of the interior of $P$ (a countable union of dyadic cubes), hence every Borel subset of the interior. The boundary of $P$ is a finite union of pieces of hyperplanes,
which are null for both measures: a piece $[0, 1]^{n-1} \times \{0\}$ is covered by $j^{n-1}$ cubes of diameter $\sqrt n/j$, with $\sum\alpha(n)(\operatorname{diam})^n \to 0$ as $j \to \infty$. So
the two measures agree on the Borel subsets of $P$, and letting $m \to \infty$, on all Borel sets.

\emph{The constant is $1$.} Let $\delta > 0$. For every cover of $Q$ by sets $C_k$, the isodiametric inequality (Theorem~\ref{geometric-measure-theory:theorem-isodiametric}) gives
$\sum_k\alpha(n)(\operatorname{diam}C_k)^n \geq \sum_k\lambda^*(C_k) \geq \lambda(Q) = 1$, so $c = \mathcal{H}^n(Q) \geq 1$. Conversely let $U = (0, 1)^n$. By the Vitali covering theorem
(Theorem~\ref{lebesgue-differentiation:theorem-vitali-covering}) there are disjoint closed balls $\overline{B}_i \subseteq U$ of diameter at most $\delta$ such that $N = U \setminus \bigcup_i\overline{B}_i$ is a
Lebesgue null set. $N$ is a Borel set, so $\mathcal{H}^n_\delta(N) \leq \mathcal{H}^n(N) = c\lambda(N) = 0$ by the comparison just proved, and
\[
\mathcal{H}^n_\delta(U) \leq \sum_i\alpha(n)(\operatorname{diam}\overline{B}_i)^n + \mathcal{H}^n_\delta(N) = \sum_i\lambda(\overline{B}_i) \leq \lambda(U) = 1.
\]
Letting $\delta \to 0$, $c = c\lambda(U) = \mathcal{H}^n(U) \leq 1$.

The last statement follows from Proposition~\ref{geometric-measure-theory:proposition-hausdorff}(4): $\mathcal{H}^n$ is finite on bounded sets and $\RR^n$ is a countable union of them, so
$\mathcal{H}^t = 0$ for $t > n$, while a nonempty open set has positive $\mathcal{H}^n$-measure.
\end{proof}

\begin{example}\label{geometric-measure-theory:example-cantor-dimension}
\emph{The Cantor set has Hausdorff dimension $s = \log 2/\log 3$}, and $\alpha(s)/2 \leq \mathcal{H}^s(C) \leq \alpha(s)$; note $3^s = 2$. Let $C \subseteq [0, 1]$ be as in
Counterexample~\ref{lebesgue-differentiation:counterexample-cantor-function}. For $k \geq 0$ and digits $a_1, \ldots, a_k \in \{0, 2\}$ let $I_{a_1\cdots a_k}$ be the closed interval with left end
$\sum_{j \leq k}a_j3^{-j}$ and length $3^{-k}$; there are $2^k$ of them, and $C$ is contained in their union $C_k$, since the digits beyond the $k$-th contribute at most $3^{-k}$.
\begin{enumerate}
\item \emph{Upper bound.} The $2^k$ intervals form a $3^{-k}$-cover of $C$ with $\sum\alpha(s)3^{-ks} = \alpha(s)2^k3^{-ks} = \alpha(s)$. So $\mathcal{H}^s(C) \leq \alpha(s)$.
\item \emph{Separation.} Two of these intervals of the same generation $k$ are at distance at least $3^{-k}$: if their digit strings first differ at position $j \leq k$, their left ends differ by
at least $2 \cdot 3^{-j} - \sum_{j < i \leq k}2 \cdot 3^{-i} = 3^{-j} + 3^{-k} \geq 2 \cdot 3^{-k}$. So a set of diameter less than $3^{-k}$ meets at most one of them.
\item \emph{The Cantor measure.} Let $\mu = \mu_F$ be the measure of Lemma~\ref{lebesgue-differentiation:lemma-stieltjes} for the Cantor function $F$. Since $F$ is constant on each interval of
$[0, 1] \setminus C$, $\mu([0, 1] \setminus C) = 0$ and $\mu(C) = F(1) - F(0) = 1$. On the generation-$k$ interval $I$ with left end $a = \sum_{j \leq k}a_j3^{-j}$ and right end
$a + \sum_{j > k}2 \cdot 3^{-j}$, $\mu(I) = F(a + 3^{-k}) - F(a) = \sum_{j > k}2^{-j} = 2^{-k}$, by the formula for $c$ on $C$.
\item \emph{Lower bound.} Let $U$ be a closed set with $0 < \operatorname{diam}U = d < 1$, and $k$ with $3^{-k-1} \leq d < 3^{-k}$. By (2), $U$ meets at most one generation-$k$ interval $I$,
and $\mu(U) \leq \mu(I) + \mu([0, 1] \setminus C) = 2^{-k} = 3^s3^{-(k+1)s} \leq 2d^s$; for a single point, $\mu(U) = 0$. For a cover of $C$ by closed sets $U_i$ of diameter less than $1$,
$1 = \mu(C) \leq \sum_i\mu(U_i \cap [0, 1]) \leq 2\sum_i(\operatorname{diam}U_i)^s$. Hence $\mathcal{H}^s_\delta(C) \geq \alpha(s)/2$ for $\delta < 1$.
\end{enumerate}
The argument in (4), bounding the mass that a measure on $A$ gives to small sets, is the \emph{mass distribution principle}, the standard way to bound Hausdorff measures from below. So the Cantor
set is too small to have positive length but too large to have zero $s$-dimensional measure for any $s < \log 2/\log 3 \approx 0.63$.
\end{example}

The estimate of Proposition~\ref{geometric-measure-theory:proposition-hausdorff}(3) for Lipschitz maps has a version for maps that are only H\"older continuous; it changes the dimension.

\begin{proposition}\label{geometric-measure-theory:proposition-holder-images}
Let $A \subseteq X$, and let $f \colon A \to Y$ be \emph{H\"older continuous} with exponent $0 < \gamma \leq 1$ and constant $L > 0$: $d(f(x), f(y)) \leq L\,d(x, y)^\gamma$ for $x, y \in A$. Then for every $s \geq 0$
\[
\mathcal{H}^{s/\gamma}(f(A)) \leq \frac{\alpha(s/\gamma)}{\alpha(s)}L^{s/\gamma}\,\mathcal{H}^s(A),
\]
and $\dim_{\mathcal{H}}f(A) \leq \dim_{\mathcal{H}}A/\gamma$.
\end{proposition}

\begin{proof}
Let $(C_k)$ be a cover of $A$ by sets of diameter at most $\delta$. The sets $f(C_k \cap A)$ cover $f(A)$ and have diameter at most $L(\operatorname{diam}C_k)^\gamma \leq L\delta^\gamma$, so, with the conventions on
empty sets and points made before Definition~\ref{geometric-measure-theory:definition-hausdorff},
\[\begin{aligned}
\sum_k\alpha(s/\gamma)(\operatorname{diam}f(C_k \cap A))^{s/\gamma} &\leq \alpha(s/\gamma)L^{s/\gamma}\sum_k(\operatorname{diam}C_k)^s \\
&= \frac{\alpha(s/\gamma)}{\alpha(s)}L^{s/\gamma}\sum_k\alpha(s)(\operatorname{diam}C_k)^s .
\end{aligned}\]
Hence $\mathcal{H}^{s/\gamma}_{L\delta^\gamma}(f(A)) \leq \frac{\alpha(s/\gamma)}{\alpha(s)}L^{s/\gamma}\mathcal{H}^s_\delta(A)$, and we let $\delta \to 0$. If $s > \dim_{\mathcal{H}}A$, then $\mathcal{H}^s(A) = 0$, so
$\mathcal{H}^{s/\gamma}(f(A)) = 0$ and $\dim_{\mathcal{H}}f(A) \leq s/\gamma$.
\end{proof}

For $\gamma = 1$ this is the Lipschitz estimate. H\"older maps can raise the dimension by the factor $1/\gamma$, and no more.

\begin{example}\label{geometric-measure-theory:example-holder-images}
\begin{enumerate}
\item \emph{Graphs and projections.} The orthogonal projection $\RR^n \to \RR^k$ onto the first coordinates is $1$-Lipschitz, so it does not increase $\mathcal{H}^s$. For a function
$g \colon [0, 1] \to \RR$, the projection maps the graph $\Gamma_g$ onto $[0, 1]$, so $\mathcal{H}^1(\Gamma_g) \geq \mathcal{H}^1([0, 1]) = 1$ by Theorem~\ref{geometric-measure-theory:theorem-hausdorff-lebesgue}: a graph is at least
as long as its domain. Similarly a set $A \subseteq \RR^2$ with $\dim_{\mathcal{H}}A < 1$ projects onto a Lebesgue null subset of every line.
\item \emph{Space-filling curves are not Lipschitz.} There are continuous surjections $g \colon [0, 1] \to [0, 1]^2$ (Theorem~\ref{general-topology:theorem-peano-curve}). If such a map were H\"older with
exponent $\gamma_0$, the proposition would give $2 = \dim_{\mathcal{H}}[0, 1]^2 \leq \dim_{\mathcal{H}}[0, 1]/\gamma_0 = 1/\gamma_0$, using Theorem~\ref{geometric-measure-theory:theorem-hausdorff-lebesgue} for both dimensions. So a space-filling curve is
never H\"older with an exponent larger than $\frac12$, in particular never Lipschitz.
\item \emph{The Cantor function is H\"older.} Let $C$ be the Cantor set and $s = \log 2/\log 3$, and let $\varphi \colon C \to [0, 1]$ send $\sum_n2a_n3^{-n}$ to $\sum_na_n2^{-n}$ ($a_n \in \{0, 1\}$), a continuous
surjection (Proposition~\ref{general-topology:proposition-cantor-surjections}). If $x \neq y$ in $C$ have digit sequences first differing at position $n$, then $|x - y| \geq 3^{-n}$
(by the proof of Proposition~\ref{general-topology:proposition-cantor-set}(1)) and $|\varphi(x) - \varphi(y)| \leq \sum_{m \geq n}2^{-m} = 2 \cdot 2^{-n} = 2(3^{-n})^s \leq 2|x - y|^s$. So $\varphi$ is H\"older with exponent $s$
and constant $2$, and the proposition with exponent $\gamma = s$, applied in dimension $s$, gives
\[
1 = \mathcal{H}^1([0, 1]) \leq \frac{\alpha(1)}{\alpha(s)}2^{1/s}\mathcal{H}^s(C) = \frac{3}{\alpha(s)}\mathcal{H}^s(C),
\]
as $\alpha(1) = 1$ and $2^{1/s} = 2^{\log 3/\log 2} = 3$. So $\mathcal{H}^s(C) \geq \alpha(s)/3 > 0$: a new proof that $\dim_{\mathcal{H}}C \geq s$, slightly weaker than the bound $\alpha(s)/2$ of
Example~\ref{geometric-measure-theory:example-cantor-dimension}, which used the Cantor measure.
\item \emph{Bi-Lipschitz invariance.} If $f$ is bi-Lipschitz, $L^{-1}d(x, y) \leq d(f(x), f(y)) \leq Ld(x, y)$, then $L^{-s}\mathcal{H}^s(A) \leq \mathcal{H}^s(f(A)) \leq L^s\mathcal{H}^s(A)$, so Hausdorff dimension, and
finiteness and positivity of $\mathcal{H}^s$, are bi-Lipschitz invariants. They are not topological invariants. On $\RR$ the \emph{snowflake metric} $d(x, y) = |x - y|^{1/2}$ is a metric (as
$\sqrt{a + b} \leq \sqrt a + \sqrt b$) with the same open sets as the usual one, since its balls are the usual balls with squared radii. A set $C$ has $d$-diameter $(\operatorname{diam}C)^{1/2}$, so the covers in the
definition give $\mathcal{H}^s_d = \frac{\alpha(s)}{\alpha(s/2)}\mathcal{H}^{s/2}$, and $[0, 1]$ has dimension $2$ for $d$. The identity $(\RR, |\cdot|) \to (\RR, d)$ is a homeomorphism, H\"older with
exponent $\frac12$, which doubles all dimensions.
\end{enumerate}
\end{example}

\section{Self-similar sets}\label{geometric-measure-theory:section-self-similar}

The Cantor set $C$ is the union of two copies of itself, $C/3$ and $C/3 + 2/3$, each scaled by the factor $\frac13$. If we only knew that $0 < \mathcal{H}^s(C) < \infty$ for some $s$, the scaling property
of Proposition~\ref{geometric-measure-theory:proposition-hausdorff}(3) would already give the value of $s$: the two copies are disjoint Borel sets, so
$\mathcal{H}^s(C) = 2 \cdot 3^{-s}\mathcal{H}^s(C)$, and $3^s = 2$. Many sets that come up in analysis and dynamics, usually called \emph{fractals}, are built in the same way from finitely many
scaled copies of themselves, and their dimension can be read off from the scaling factors. This section makes the heuristic rigorous: we construct such sets in general
(Theorem~\ref{geometric-measure-theory:theorem-hutchinson}), and compute their dimension when the copies are disjoint (Theorem~\ref{geometric-measure-theory:theorem-moran}). Along the way we meet a second notion of
dimension, which is cruder than the Hausdorff dimension but easier to compute. References: \cite{falconer-fractal}, \cite{hutchinson-1981}.

\subsection{Box-counting dimension}\label{geometric-measure-theory:subsection-box-dimension}

Hausdorff measure allows covers by sets of very different sizes, and this is what makes it a measure. If instead we cover a set by sets of one fixed size $\delta$ and only count them, we get a
simpler quantity. A segment of length $1$ needs about $1/\delta$ sets of diameter $\delta$, a square about $\delta^{-2}$, a cube about $\delta^{-3}$: the exponent is the dimension.

\begin{definition}\label{geometric-measure-theory:definition-box-dimension}
For a nonempty bounded $A \subseteq \RR^n$ and $\delta > 0$, let $N_\delta(A)$ be the smallest number of sets of diameter at most $\delta$ that cover $A$. The \emph{lower} and \emph{upper box-counting
dimensions} of $A$ are
\[
\underline{\dim}_BA = \liminf_{\delta \to 0}\frac{\log N_\delta(A)}{-\log\delta}, \qquad \overline{\dim}_BA = \limsup_{\delta \to 0}\frac{\log N_\delta(A)}{-\log\delta}.
\]
If they are equal, their common value is the \emph{box-counting dimension} $\dim_BA$.
\end{definition}

The number $N_\delta(A)$ is finite: $A$ lies in a cube, which is a finite union of small cubes. The definition says that $N_\delta(A)$ is roughly $\delta^{-\dim_BA}$; more precisely,
$\overline{\dim}_BA < t$ if and only if $N_\delta(A) \leq \delta^{-t^{\prime}}$ for some $t^{\prime} < t$ and all small $\delta$.

\begin{proposition}\label{geometric-measure-theory:proposition-dimension-properties}
Let $A, A_1, A_2, \ldots$ be subsets of $\RR^n$.
\begin{enumerate}
\item (\emph{Countable stability}) $\dim_{\mathcal{H}}\bigcup_kA_k = \sup_k\dim_{\mathcal{H}}A_k$. In particular countable sets have Hausdorff dimension $0$.
\item If $A$ is nonempty and bounded, $N_\delta(A) = N_\delta(\overline A)$, so $A$ and its closure have the same box-counting dimensions.
\item If $A$ is nonempty and bounded, $\dim_{\mathcal{H}}A \leq \underline{\dim}_BA \leq \overline{\dim}_BA \leq n$.
\end{enumerate}
\end{proposition}

\begin{proof}
(1) Since $\mathcal{H}^s$ is monotone, $\dim_{\mathcal{H}}A_k \leq \dim_{\mathcal{H}}\bigcup_kA_k$ for every $k$. Conversely, if $s > \sup_k\dim_{\mathcal{H}}A_k$, then $\mathcal{H}^s(A_k) = 0$ for every $k$, and
$\mathcal{H}^s(\bigcup_kA_k) = 0$ by countable subadditivity; so $\dim_{\mathcal{H}}\bigcup_kA_k \leq s$. A point has $\mathcal{H}^s$-measure $0$ for every $s > 0$, since it is covered by one set of
diameter $0$; so a point has dimension $0$, and so has a countable set.

(2) A cover of $A$ by sets of diameter at most $\delta$ gives the cover of $\overline A$ by their closures, which have the same diameters, since a finite union of closed sets is closed. So
$N_\delta(\overline A) \leq N_\delta(A)$, and the other inequality holds as $A \subseteq \overline A$.

(3) The middle inequality is clear. For the last one, let $A$ lie in a cube of side $\ell$. Dividing it into $m^n$ cubes of side $\ell/m$, of diameter $\ell\sqrt n/m$, with $m$ the smallest integer
$\geq \ell\sqrt n/\delta$, gives $N_\delta(A) \leq m^n \leq (\ell\sqrt n/\delta + 1)^n$, and $\log N_\delta(A)/(-\log\delta) \leq n\log(\ell\sqrt n + \delta)/(-\log\delta) + n$, whose $\limsup$ as
$\delta \to 0$ is $n$. For the first inequality, let $s > t > \underline{\dim}_BA$. There are $\delta_j \to 0$ with $\log N_{\delta_j}(A) < t(-\log\delta_j)$, that is, $N_{\delta_j}(A) < \delta_j^{-t}$.
Covering $A$ by $N_{\delta_j}(A)$ sets of diameter at most $\delta_j$ gives
\[
\mathcal{H}^s_{\delta_j}(A) \leq \alpha(s)N_{\delta_j}(A)\delta_j^s \leq \alpha(s)\delta_j^{s - t} \to 0 .
\]
As $\mathcal{H}^s_\delta(A)$ increases when $\delta$ decreases, $\mathcal{H}^s(A) = \lim_j\mathcal{H}^s_{\delta_j}(A) = 0$, so $\dim_{\mathcal{H}}A \leq s$. Since $s > \underline{\dim}_BA$ was arbitrary, the
inequality follows.
\end{proof}

By (2), box-counting dimension cannot see the difference between a set and its closure: $\QQ \cap [0, 1]$ has box-counting dimension $1$, like $[0, 1]$, while its Hausdorff dimension is $0$ by (1).
So box-counting dimension is not countably stable, and in general it is not attached to a measure. The inequality in (3) can be strict even for compact sets.

\begin{example}\label{geometric-measure-theory:example-box-hausdorff}
\emph{A compact countable set of box-counting dimension $\frac12$.} Let $E = \{0\} \cup \{1/m : m \geq 1\}$. It is countable, so $\dim_{\mathcal{H}}E = 0$. We show that $\dim_BE = \frac12$. Let
$0 < \delta < 1$.
\begin{enumerate}
\item \emph{Lower bound.} Let $m = \lfloor\delta^{-1/2}\rfloor \geq 1$. For $2 \leq j \leq m$ the consecutive points $1/(j-1)$ and $1/j$ are at distance
$\frac1{j(j-1)} > \frac1{m^2} \geq \delta$. So the points $1, \frac12, \ldots, \frac1m$ are pairwise more than $\delta$ apart, a set of diameter at most $\delta$ contains at most one of them, and
$N_\delta(E) \geq m > \delta^{-1/2} - 1$.
\item \emph{Upper bound.} Let $m = \lceil\delta^{-1/2}\rceil$. The interval $[0, \frac1m]$, which contains all points $1/j$ with $j \geq m$, is covered by $\lceil\frac1{m\delta}\rceil$ intervals of
length $\delta$, and the remaining $m - 1$ points are covered by singletons. Since $m - 1 < \delta^{-1/2}$ and $\frac1{m\delta} \leq \delta^{-1/2}$, $N_\delta(E) \leq 2\delta^{-1/2} + 1 \leq 3\delta^{-1/2}$.
\end{enumerate}
Taking logarithms, $\log(\delta^{-1/2} - 1)/(-\log\delta)$ and $(\log 3 + \frac12(-\log\delta))/(-\log\delta)$ both tend to $\frac12$ as $\delta \to 0$. The reason is that the points $1/j$ are spread
out for $j$ small and crowded for $j$ large: at scale $\delta$ one sees about $\delta^{-1/2}$ separate points, and the rest look like an interval of length about $\delta^{1/2}$.
\end{example}

\subsection{Iterated function systems}\label{geometric-measure-theory:subsection-iterated-function-systems}

A map $f \colon \RR^n \to \RR^n$ is a \emph{similarity with ratio $r > 0$} if $|f(x) - f(y)| = r|x - y|$ for all $x, y$; for example $f(x) = rOx + b$ with $O$ orthogonal (every similarity has this
form, but we do not need this). A similarity with ratio $r$ multiplies diameters by $r$, and by Proposition~\ref{geometric-measure-theory:proposition-hausdorff}(3), applied to $f$ and to its inverse,
$\mathcal{H}^s(f(A)) = r^s\mathcal{H}^s(A)$ for every $A$ and $s$.

\begin{definition}\label{geometric-measure-theory:definition-ifs}
An \emph{iterated function system} (of similarities) is a finite family $f_1, \ldots, f_N$ of similarities of $\RR^n$ with ratios $r_1, \ldots, r_N \in (0, 1)$, where $N \geq 2$. For a set $A$ write
$F(A) = f_1(A) \cup \cdots \cup f_N(A)$. A nonempty compact set $K$ with $F(K) = K$ is called the \emph{attractor} of the system, or a \emph{self-similar set}.
\end{definition}

It is convenient to use words in the alphabet $\{1, \ldots, N\}$. For a word $w = w_1 \cdots w_k$ of length $|w| = k$ let
\[
f_w = f_{w_1} \circ f_{w_2} \circ \cdots \circ f_{w_k}, \qquad r_w = r_{w_1}r_{w_2} \cdots r_{w_k},
\]
so that $f_w$ is a similarity with ratio $r_w$; for the empty word $\varnothing$, $f_\varnothing$ is the identity and $r_\varnothing = 1$. Write $wi$ for the word $w$ followed by the letter $i$, and
$r = \max_ir_i < 1$, so that $r_w \leq r^{|w|}$. For an infinite sequence $\omega = \omega_1\omega_2\cdots$ of letters, $\omega|k = \omega_1 \cdots \omega_k$ is its beginning of length $k$. Since
$F(A) = \bigcup_if_i(A)$, induction on $k$ gives $F^k(A) = \bigcup_{|w| = k}f_w(A)$.

\begin{theorem}[Hutchinson]\label{geometric-measure-theory:theorem-hutchinson}
Every iterated function system has a unique attractor $K$. Moreover:
\begin{enumerate}
\item If $B$ is a nonempty compact set with $F(B) \subseteq B$, then $K = \bigcap_kF^k(B)$; in particular $K \subseteq B$.
\item Let $K_w = f_w(K)$. Then $K = \bigcup_{|w| = k}K_w$ for every $k$, and $K_w = \bigcup_iK_{wi}$, with $\operatorname{diam}K_w = r_w\operatorname{diam}K$.
\item Every $x \in K$ has an \emph{address}: an infinite sequence $\omega$ with $x \in K_{\omega|k}$ for all $k$. For every infinite sequence $\omega$, $\bigcap_kK_{\omega|k}$ is a single point.
\end{enumerate}
\end{theorem}

\begin{proof}
\emph{Existence, and (1).} First, a compact set $B$ with $F(B) \subseteq B$ exists: let $R = \max_i|f_i(0)|/(1 - r)$ and let $B$ be the closed ball of radius $R$ about $0$. For $|x| \leq R$,
$|f_i(x)| \leq |f_i(x) - f_i(0)| + |f_i(0)| \leq rR + (1 - r)R = R$. Now let $B$ be any nonempty compact set with $F(B) \subseteq B$. Applying $F$, which preserves inclusions, gives
$B \supseteq F(B) \supseteq F^2(B) \supseteq \cdots$, a decreasing sequence of nonempty compact sets (finite unions of continuous images of compact sets). Their intersection $K$ is nonempty and compact. To
see $F(K) = K$, note that $f_i(K) = \bigcap_kf_i(F^k(B))$ since $f_i$ is injective, so
\[
F(K) = \bigcup_i\bigcap_kf_i(F^k(B)) \qquad\text{and}\qquad K = \bigcap_kF^{k+1}(B) = \bigcap_k\bigcup_if_i(F^k(B)).
\]
The first set is contained in the second. Conversely, if $x$ lies in $\bigcup_if_i(F^k(B))$ for every $k$, then there is one $i$ with $x \in f_i(F^k(B))$ for infinitely many $k$, and so for all $k$, as
these sets decrease with $k$. So the two sets are equal.

\emph{Uniqueness.} Let $A$ and $K$ be nonempty compact sets with $F(A) = A$ and $F(K) = K$, and $D = \max\{|a - y| : a \in A, y \in K\}$. Fix $k$. Since $A = F^k(A)$, every $a \in A$ is $f_w(a^{\prime})$
for a word $w$ of length $k$ and some $a^{\prime} \in A$. For any $y \in K$, the point $f_w(y)$ lies in $F^k(K) = K$, and $|a - f_w(y)| = r_w|a^{\prime} - y| \leq r^kD$. So every point of $A$ is
within $r^kD$ of $K$, for every $k$; as $K$ is closed, $A \subseteq K$. By symmetry $K \subseteq A$.

(2) The first statement is $K = F^k(K)$, the second is $f_w(K) = f_w(F(K)) = \bigcup_if_w(f_i(K))$, and the third holds because $f_w$ multiplies distances by $r_w$.

(3) Since $K = \bigcup_if_i(K)$, a point $x \in K$ can be written $x = f_{\omega_1}(x_1)$ with $x_1 \in K$; then $x_1 = f_{\omega_2}(x_2)$ with $x_2 \in K$, and so on. Then
$x = f_{\omega|k}(x_k) \in K_{\omega|k}$ for all $k$. For any $\omega$, the sets $K_{\omega|k}$ are nonempty, compact, decreasing by (2), and of diameter at most $r^k\operatorname{diam}K \to 0$; so
their intersection is one point.
\end{proof}

\begin{remark}\label{geometric-measure-theory:remark-hutchinson}
Hutchinson's original proof \cite{hutchinson-1981} is an application of the Banach fixed point theorem. The nonempty compact subsets of $\RR^n$ form a complete metric space for the
\emph{Hausdorff distance} $d_H(A, B) = \inf\{\varepsilon > 0 : A \subseteq B_\varepsilon, B \subseteq A_\varepsilon\}$, where $A_\varepsilon$ is the set of points at distance less than $\varepsilon$ from
$A$, and $F$ is a contraction with constant $r$ for this distance; the attractor is its fixed point, and $F^k(B) \to K$ for every nonempty compact $B$. This also explains the name: starting from any
picture $B$, repeatedly replacing it by $N$ smaller copies of itself converges to $K$. The uniqueness proof above is the contraction argument in disguise.
\end{remark}

\begin{example}\label{geometric-measure-theory:example-ifs}
\begin{enumerate}
\item \emph{The Cantor set.} $f_1(x) = x/3$, $f_2(x) = x/3 + 2/3$ on $\RR$. The Cantor set $C$, the set of $\sum_ka_k3^{-k}$ with all $a_k \in \{0, 2\}$, satisfies $f_1(C) = \{a_1 = 0\}$ and
$f_2(C) = \{a_1 = 2\}$, so $F(C) = C$ and $C$ is the attractor. The sets $C_w$ are the parts of $C$ in the intervals $I_w$ of Example~\ref{geometric-measure-theory:example-cantor-dimension}, and the
address of a point is its sequence of ternary digits, with $1$ and $2$ standing for the digits $0$ and $2$.
\item \emph{An interval.} $f_1(x) = x/2$, $f_2(x) = x/2 + 1/2$. Here $F([0, 1]) = [0, 1]$, so the attractor is $[0, 1]$; addresses are binary expansions, and the dyadic rationals have two addresses.
\item \emph{The Sierpi\'nski gasket.} Let $p_1, p_2, p_3$ be the vertices of a triangle $T$ in $\RR^2$ and $f_i(x) = (x + p_i)/2$. Then $f_i(T)$ is the triangle with vertex $p_i$ and the midpoints of the two
sides through $p_i$, so $F(T) \subseteq T$ and the attractor is $S = \bigcap_kF^k(T)$ by (1): at each step every triangle is replaced by its three corner triangles, and the open middle one is
removed.
\item \emph{The Koch curve.} In $\RR^2 = \CC$ let $\zeta = e^{i\pi/3}$ and
\[
f_1(z) = \tfrac z3, \quad f_2(z) = \tfrac{\zeta z + 1}{3}, \quad f_3(z) = \tfrac{\bar{\zeta}z}{3} + \tfrac12 + \tfrac{i\sqrt3}{6}, \quad f_4(z) = \tfrac{z + 2}{3}.
\]
They map the segment $[0, 1]$ to the four sides of the polygon $0, \frac13, \frac12 + \frac{i\sqrt3}{6}, \frac23, 1$, whose middle part is the tip of an equilateral triangle erected on the middle third
of $[0, 1]$: for instance $f_2(1) = \frac{1 + \zeta}{3} = \frac12 + \frac{i\sqrt3}{6} = f_3(0)$ and $f_3(1) = \frac{\bar{\zeta}}{3} + \frac12 + \frac{i\sqrt3}{6} = \frac23 = f_4(0)$. The attractor is the
limit of the polygons $F^k([0, 1])$, in the sense of Remark~\ref{geometric-measure-theory:remark-hutchinson}.
\item \emph{The four-corner Cantor set.} For $a \in \{0, \frac34\}^2$ let $f_a(x) = \frac x4 + a$ on $\RR^2$. The four squares $f_a([0, 1]^2)$ are the corner squares of side $\frac14$ of the unit
square, so the attractor $K_4 \subseteq [0, 1]^2$ is obtained by repeatedly replacing every square by its four corner squares of a quarter of the size. It contains the fixed points $(0, 0)$ and
$(1, 1)$ of $f_{(0, 0)}$ and $f_{(3/4, 3/4)}$, so $\operatorname{diam}K_4 = \sqrt2$.
\end{enumerate}
\end{example}

\subsection{The similarity dimension}\label{geometric-measure-theory:subsection-similarity-dimension}

The function $s \mapsto \sum_ir_i^s$ is continuous and strictly decreasing on $[0, \infty)$, equal to $N \geq 2$ at $s = 0$ and tending to $0$ as $s \to \infty$. So there is a unique $s > 0$ with
\begin{equation}\label{geometric-measure-theory:equation-moran}
r_1^s + r_2^s + \cdots + r_N^s = 1,
\end{equation}
the \emph{similarity dimension} of the system. If all ratios equal $\rho$, then $N\rho^s = 1$ and $s = \log N/\log(1/\rho)$: for the Cantor set $s = \log 2/\log 3$, for the interval of
Example~\ref{geometric-measure-theory:example-ifs}(2) $s = 1$, for the Sierpi\'nski gasket $s = \log 3/\log 2 \approx 1.585$, for the Koch curve $s = \log 4/\log 3 \approx 1.262$, and for the
four-corner set $s = \log 4/\log 4 = 1$. To cover $K$ efficiently by the pieces $K_w$ one should use pieces of comparable sizes, not of equal word length; the following lemma organises this.

\begin{lemma}\label{geometric-measure-theory:lemma-cut-sets}
For $0 < \rho < 1$ let $Q_\rho$ be the set of nonempty words $w$ with $r_w \leq \rho < r_{w^-}$, where $w^-$ is $w$ with its last letter removed. Then every infinite sequence $\omega$ has exactly one
beginning $\omega|k$ in $Q_\rho$, the set $Q_\rho$ is finite, and
\[
\sum_{w \in Q_\rho}r_w^s = 1, \qquad \#Q_\rho \leq (\rho\,r_{\min})^{-s}, \qquad K = \bigcup_{w \in Q_\rho}K_w,
\]
where $s$ is the similarity dimension and $r_{\min} = \min_ir_i$.
\end{lemma}

\begin{proof}
Along $\omega$ the numbers $r_{\omega|k}$ decrease strictly from $r_{\omega|0} = 1 > \rho$ to $0$, so there is exactly one $k \geq 1$ with $r_{\omega|k} \leq \rho < r_{\omega|(k - 1)}$. Let $k_0$ be such
that $r^{k_0} \leq \rho$. Every word $w$ of length at least $k_0 + 1$ has $r_{w^-} \leq r^{k_0} \leq \rho$, so it is not in $Q_\rho$; hence $Q_\rho$ is finite. Every word $v$ of length $k_0$ has
$r_v \leq \rho$, so, applying the first statement to any infinite sequence beginning with $v$, $v$ has exactly one beginning in $Q_\rho$. Grouping the words of length $k_0$ according to this beginning,
and using that the words of length $k_0$ extending a given $w$ are the $wu$ with $|u| = k_0 - |w|$,
\[
1 = \Bigl(\sum_ir_i^s\Bigr)^{k_0} = \sum_{|v| = k_0}r_v^s = \sum_{w \in Q_\rho}\ \sum_{|u| = k_0 - |w|}r_w^sr_u^s = \sum_{w \in Q_\rho}r_w^s\Bigl(\sum_ir_i^s\Bigr)^{k_0 - |w|} = \sum_{w \in Q_\rho}r_w^s .
\]
Each $w \in Q_\rho$ has $r_w = r_{w^-}r_{w_{|w|}} > \rho\,r_{\min}$, so $1 \geq \#Q_\rho\,(\rho\,r_{\min})^s$. Finally, if $x \in K$ has address $\omega$ and $\omega|k \in Q_\rho$, then
$x \in K_{\omega|k}$.
\end{proof}

\begin{proposition}\label{geometric-measure-theory:proposition-similarity-upper}
Let $K$ be the attractor of an iterated function system with similarity dimension $s$. Then $\mathcal{H}^s(K) \leq \alpha(s)(\operatorname{diam}K)^s$, and $\dim_{\mathcal{H}}K \leq \overline{\dim}_BK \leq s$.
\end{proposition}

\begin{proof}
We may assume $d = \operatorname{diam}K > 0$. For $0 < \rho < 1$ the sets $K_w$, $w \in Q_\rho$, cover $K$ by Lemma~\ref{geometric-measure-theory:lemma-cut-sets} and have diameter
$r_wd \leq \rho d$, so
\[
\mathcal{H}^s_{\rho d}(K) \leq \sum_{w \in Q_\rho}\alpha(s)(r_wd)^s = \alpha(s)d^s,
\]
and $\rho \to 0$ gives the first claim. With $\delta = \rho d$ they also show $N_\delta(K) \leq \#Q_\rho \leq (r_{\min}\delta/d)^{-s}$ for $0 < \delta < d$, so
$\log N_\delta(K)/(-\log\delta) \leq s + s\log(d/r_{\min})/(-\log\delta)$, whose $\limsup$ is $s$. The remaining inequality is
Proposition~\ref{geometric-measure-theory:proposition-dimension-properties}(3).
\end{proof}

The upper bound can be far from the truth when the pieces overlap: for $f_1(x) = f_2(x) = x/2$ the attractor is the point $\{0\}$, while $s = 1$. The lower bound needs a condition which prevents such
overlaps. The simplest one is that the pieces are disjoint.

\begin{theorem}\label{geometric-measure-theory:theorem-moran}
Let $K$ be the attractor of an iterated function system with similarity dimension $s$, and suppose that the sets $f_1(K), \ldots, f_N(K)$ are pairwise disjoint (the \emph{strong separation condition}).
Then
\[
\alpha(s)(r_{\min}\eta)^s \leq \mathcal{H}^s(K) \leq \alpha(s)(\operatorname{diam}K)^s, \qquad \eta = \min_{i \neq j}d(f_i(K), f_j(K)) > 0,
\]
and $\dim_{\mathcal{H}}K = \dim_BK = s$.
\end{theorem}

\begin{proof}
\emph{Step 1: separation at all scales.} $\eta > 0$ because the $f_i(K)$ are disjoint compact sets. For a word $w$ and $i \neq j$, $f_w$ multiplies distances by $r_w$, so
$d(K_{wi}, K_{wj}) = r_wd(f_i(K), f_j(K)) \geq r_w\eta$. Let $x, y \in K$ with addresses $\omega$ and $\omega^{\prime}$, and suppose these first differ at position $j$; let $u = \omega|(j - 1)$. Then
$x \in K_{u\omega_j}$ and $y \in K_{u\omega^{\prime}_j}$, so
\begin{equation}\label{geometric-measure-theory:equation-separation}
|x - y| \geq r_{\omega|(j - 1)}\,\eta .
\end{equation}
In particular, for each $k$ the $N^k$ sets $K_w$ with $|w| = k$ are pairwise disjoint, and each of them is open in $K$, its complement in $K$ being the finite union of the others, which is compact.

\emph{Step 2: the natural measure.} Put $p_w = r_w^s$, so that $\sum_ip_{wi} = p_w\sum_ir_i^s = p_w$, and by induction $\sum_{|u| = l}p_{wu} = p_w$ for every $l$. Fix $x_0 \in K$, put
$x_w = f_w(x_0) \in K_w$, and for $\varphi \in C(K)$ let
\[
S_k(\varphi) = \sum_{|w| = k}p_w\,\varphi(x_w) .
\]
Let $\omega_\varphi(t) = \sup\{|\varphi(x) - \varphi(y)| : |x - y| \leq t\}$, which tends to $0$ as $t \to 0$ since $\varphi$ is uniformly continuous on the compact set $K$. For $l > k$, regrouping the words of
length $l$ by their beginnings of length $k$,
\[\begin{aligned}
|S_l(\varphi) - S_k(\varphi)| &= \Bigl|\sum_{|w| = k}\sum_{|u| = l - k}p_{wu}\bigl(\varphi(x_{wu}) - \varphi(x_w)\bigr)\Bigr| \\
&\leq \omega_\varphi(r^k\operatorname{diam}K)\sum_{|w| = k}p_w = \omega_\varphi(r^k\operatorname{diam}K),
\end{aligned}\]
since $x_{wu}$ and $x_w$ both lie in $K_w$. So $(S_k(\varphi))_k$ is a Cauchy sequence, and $\Lambda(\varphi) = \lim_kS_k(\varphi)$ is a positive linear functional on $C(K) = C_c(K)$ with $\Lambda(1) = 1$. By the
Riesz--Markov theorem (Theorem~\ref{radon-measures:theorem-riesz-markov}) there is a Borel probability measure $\mu$ on $K$ with $\int\varphi\,d\mu = \Lambda(\varphi)$. By Step 1 the indicator
function $\mathbf{1}_{K_w}$ is continuous on $K$, and for $l \geq k = |w|$ the points $x_v$ with $|v| = l$ lying in $K_w$ are those with $v$ beginning with $w$. Hence
\[
\mu(K_w) = \lim_lS_l(\mathbf{1}_{K_w}) = \sum_{|u| = l - k}p_{wu} = p_w = r_w^s .
\]

\emph{Step 3: the mass of small sets.} Let $U$ be a closed set with $d = \operatorname{diam}U < \eta$; we show $\mu(U \cap K) \leq (d/(r_{\min}\eta))^s$. This is clear if $U \cap K = \emptyset$. Let
$x \in U \cap K$ have address $\omega$. If $d = 0$, then $\mu(U \cap K) = \mu(\{x\}) \leq \mu(K_{\omega|k}) = r_{\omega|k}^s \to 0$. If $d > 0$, the numbers $r_{\omega|k}\eta$ decrease from $\eta > d$ to $0$,
so there is a least $k \geq 1$ with $r_{\omega|k}\eta \leq d$, and $r_{\omega|(k - 1)}\eta > d$. Every $y \in U \cap K$ lies in $K_{\omega|(k - 1)}$: otherwise an address of $y$ first differs from
$\omega$ at a position $j \leq k - 1$, and \eqref{geometric-measure-theory:equation-separation} gives $|x - y| \geq r_{\omega|(j - 1)}\eta \geq r_{\omega|(k - 1)}\eta > d$, which is impossible. Therefore
\[
\mu(U \cap K) \leq \mu(K_{\omega|(k - 1)}) = r_{\omega|(k - 1)}^s \leq \Bigl(\frac{d}{r_{\min}\eta}\Bigr)^s,
\]
the last step because $d \geq r_{\omega|k}\eta = r_{\omega|(k - 1)}r_{\omega_k}\eta \geq r_{\omega|(k - 1)}r_{\min}\eta$.

\emph{Step 4: conclusion.} Let $\delta < \eta$ and let $(U_i)$ be a cover of $K$ by closed sets of diameter at most $\delta$; closed covers suffice to compute $\mathcal{H}^s_\delta$. By Step 3,
\[
1 = \mu(K) \leq \sum_i\mu(U_i \cap K) \leq (r_{\min}\eta)^{-s}\sum_i(\operatorname{diam}U_i)^s,
\]
so $\mathcal{H}^s_\delta(K) \geq \alpha(s)(r_{\min}\eta)^s$, and the lower bound follows. With Proposition~\ref{geometric-measure-theory:proposition-similarity-upper}, $0 < \mathcal{H}^s(K) < \infty$, so
$\dim_{\mathcal{H}}K = s$, and $s = \dim_{\mathcal{H}}K \leq \underline{\dim}_BK \leq \overline{\dim}_BK \leq s$ by Proposition~\ref{geometric-measure-theory:proposition-dimension-properties}(3).
\end{proof}

Step 3 is the mass distribution principle of Example~\ref{geometric-measure-theory:example-cantor-dimension} again; the new ingredient is that the measure is built for us by the Riesz--Markov theorem, once we
know what it should give to the pieces $K_w$. The measure $\mu$ is in fact the normalised restriction of $\mathcal{H}^s$ to $K$ (Exercise~\ref{geometric-measure-theory:exercise-natural-measure}).

\begin{example}\label{geometric-measure-theory:example-similarity-dimension}
\begin{enumerate}
\item \emph{The Cantor set.} $f_1(C) \subseteq [0, \frac13]$ and $f_2(C) \subseteq [\frac23, 1]$ are disjoint, at distance $\eta = \frac13$, the distance between the points $\frac13$ and $\frac23$ of $C$.
The theorem gives $\dim_{\mathcal{H}}C = \dim_BC = \log 2/\log 3$ and $\alpha(s)/4 \leq \mathcal{H}^s(C) \leq \alpha(s)$, as $(r_{\min}\eta)^s = 9^{-s} = \frac14$; the direct argument of
Example~\ref{geometric-measure-theory:example-cantor-dimension} gave the better lower bound $\alpha(s)/2$. In fact $\mathcal{H}^s(C) = \alpha(s)$ \cite{falconer-fractal}, but this requires a finer argument.
\item \emph{A totally disconnected set of dimension $1$.} For the four-corner set $K_4$ of Example~\ref{geometric-measure-theory:example-ifs}(5) the four pieces lie in the four corner squares of side $\frac14$,
which are pairwise at distance at least $\frac12$. So $\dim_{\mathcal{H}}K_4 = 1$ and $0 < \mathcal{H}^1(K_4) \leq \alpha(1)\sqrt2 = \sqrt2$ (note $\alpha(1) = 1$). Yet $K_4$ contains no segment, indeed no
connected set with more than one point: the sets $(K_4)_w$ are open and closed in $K_4$ by Step 1, and every point has such neighbourhoods of arbitrarily small diameter. So $K_4$ has the dimension and
the finite positive length of a curve, but none of its geometry; Example~\ref{geometric-measure-theory:example-four-corner-projections} and Section~\ref{geometric-measure-theory:section-densities} return to this.
\item \emph{Products.} $C \times C \subseteq \RR^2$ is the attractor of the four maps $(x, y) \mapsto (x/3 + a, y/3 + b)$ with $a, b \in \{0, \frac23\}$ (by uniqueness, since $C \times C$ is compact and
invariant), whose images lie in the four corner squares of side $\frac13$. So $\dim_{\mathcal{H}}(C \times C) = \log 4/\log 3 = 2\dim_{\mathcal{H}}C$.
\item \emph{Middle-$\lambda$ Cantor sets.} Removing open middle intervals of relative length $\lambda \in (0, 1)$ instead of $\frac13$ gives the attractor of $x \mapsto \rho x$ and
$x \mapsto \rho x + 1 - \rho$ with $\rho = (1 - \lambda)/2$, of dimension $\log 2/\log(2/(1 - \lambda))$. As $\lambda$ runs through $(0, 1)$ this takes every value in $(0, 1)$.
\end{enumerate}
\end{example}

\begin{example}\label{geometric-measure-theory:example-four-corner-projections}
\emph{Shadows of the four-corner set.} Let $K_4$ be as in Example~\ref{geometric-measure-theory:example-ifs}(5). Its projection $\pi_1(x, y) = x$ to the first axis satisfies $\pi_1 \circ f_a = g_{a_1} \circ \pi_1$ with
$g_0(t) = t/4$ and $g_{3/4}(t) = t/4 + \frac34$. So $\pi_1(K_4)$ is a nonempty compact set with $\pi_1(K_4) = g_0(\pi_1(K_4)) \cup g_{3/4}(\pi_1(K_4))$, the attractor of $\{g_0, g_{3/4}\}$, whose pieces lie in
$[0, \frac14]$ and $[\frac34, 1]$. By Theorem~\ref{geometric-measure-theory:theorem-moran} it has dimension $\log 2/\log 4 = \frac12$, hence length $0$. Now project instead by $\pi(x, y) = x + y/2$. Then
$\pi \circ f_a = h_a \circ \pi$ with $h_a(t) = t/4 + \pi(a)$, and $\pi(a)$ runs through $0, \frac38, \frac34, \frac98$ as $a$ runs through $(0, 0), (0, \frac34), (\frac34, 0), (\frac34, \frac34)$. The interval
$[0, \frac32]$ satisfies
\[
\bigcup_ah_a\bigl([0, \tfrac32]\bigr) = [0, \tfrac38] \cup [\tfrac38, \tfrac34] \cup [\tfrac34, \tfrac98] \cup [\tfrac98, \tfrac32] = [0, \tfrac32],
\]
so by uniqueness of attractors $\pi(K_4) = [0, \frac32]$. The orthogonal projection onto the line spanned by $(2, 1)$ is $\frac{2}{\sqrt5}\pi$ up to an isometry of the line, so the shadow of $K_4$ in this
direction has length $\frac{3}{\sqrt5}$, while in the direction of the first axis it has length $0$. As orthogonal projections are $1$-Lipschitz, this also gives
\[
\frac{3}{\sqrt5} \leq \mathcal{H}^1(K_4) \leq \sqrt2,
\]
that is, $1.341\ldots \leq \mathcal{H}^1(K_4) \leq 1.414\ldots$, a much better lower bound than Theorem~\ref{geometric-measure-theory:theorem-moran} gives.
\end{example}

\begin{remark}\label{geometric-measure-theory:remark-open-set-condition}
\emph{Touching pieces.} For the Sierpi\'nski gasket and the Koch curve the pieces $f_i(K)$ touch, so Theorem~\ref{geometric-measure-theory:theorem-moran} does not apply, and
Proposition~\ref{geometric-measure-theory:proposition-similarity-upper} only gives the upper bounds $\log 3/\log 2$ and $\log 4/\log 3$ for their dimensions. Equality holds under
the weaker \emph{open set condition}: there is a nonempty bounded open set $V$ with $f_i(V) \subseteq V$ and the $f_i(V)$ pairwise disjoint (for the gasket, the open triangle; for the Koch curve, the
open isosceles triangle with base $[0, 1]$ and apex $\frac12 + \frac{i\sqrt3}{6}$). Under this condition $0 < \mathcal{H}^s(K) < \infty$ and $\dim_{\mathcal{H}}K = \dim_BK = s$
\cite{hutchinson-1981}, \cite[Chapter 9]{falconer-fractal}. The proof follows the same plan: the natural measure exists without any condition, and the open set condition replaces
\eqref{geometric-measure-theory:equation-separation} by a volume argument showing that a ball of radius $\rho$ meets a bounded number of pieces $K_w$, $w \in Q_\rho$.
\end{remark}

\section{Lipschitz maps and Rademacher's theorem}\label{geometric-measure-theory:section-rademacher}

A Lipschitz function of one variable is absolutely continuous, hence differentiable almost everywhere (Example~\ref{lebesgue-differentiation:example-absolutely-continuous}). Rademacher's theorem extends
this to several variables. It is the reason why Lipschitz maps are the natural class of maps in geometric measure theory: they are just regular enough to have a derivative almost everywhere,
and they are preserved by the operations of the theory (distance functions, compositions, restrictions, inverses of bi-Lipschitz maps), which smooth maps are not.

\begin{lemma}\label{geometric-measure-theory:lemma-vanishing-integrals}
If $g \in L^1_{\mathrm{loc}}(\RR^n)$ and $\int g\varphi\,d\lambda = 0$ for all compactly supported functions $\varphi$ of class $C^1$, then $g = 0$ almost everywhere.
\end{lemma}

\begin{proof}
Let $\psi(s) = 3s^2 - 2s^3$ for $0 \leq s \leq 1$, $\psi = 0$ for $s < 0$ and $\psi = 1$ for $s > 1$, a $C^1$ function. For a ball $B(x, r)$ and $j > 1/r$ let
$\varphi_j(y) = \psi(j(r - \|y - x\|))$; it is $C^1$, since it is constant near $y = x$, takes values in $[0, 1]$, vanishes outside $B(x, r)$, and $\varphi_j \to \mathbf{1}_{B(x, r)}$
pointwise. By dominated convergence $\int_{B(x, r)}g = 0$ for all balls, and by the Lebesgue differentiation theorem (Theorem~\ref{lebesgue-differentiation:theorem-lebesgue-differentiation}),
$g(x) = \lim_{r \to 0}A_rg(x) = 0$ for almost every $x$.
\end{proof}

\begin{theorem}[Rademacher]\label{geometric-measure-theory:theorem-rademacher}
Every Lipschitz map $f \colon \RR^n \to \RR^m$ is differentiable at almost every point of $\RR^n$.
\end{theorem}

\begin{proof}
It suffices to treat $m = 1$, applying the result to the components. Let $L$ be a Lipschitz constant. For a unit vector $v$ and $x \in \RR^n$ let
$D_vf(x) = \lim_{t \to 0}(f(x + tv) - f(x))/t$ where the limit exists.

\emph{Step 1: for each $v$, $D_vf(x)$ exists for almost every $x$.} As the difference quotients are continuous in $t \neq 0$, the limit exists if and only if the $\limsup$ and $\liminf$ over
rational $t \to 0$ are equal and finite, so the set $N_v$ where it fails is Borel, and $D_vf$ is Borel measurable where defined, with $|D_vf| \leq L$. Choose an orthogonal map $R$ with $Re_1 = v$;
it preserves $\lambda$ (Lemma~\ref{measure-theory:lemma-linear-change}), and $f \circ R$ is Lipschitz. For each $y \in \RR^{n-1}$, $t \mapsto f(R(t, y))$ is a Lipschitz function of one variable,
differentiable almost everywhere (Example~\ref{lebesgue-differentiation:example-absolutely-continuous}); so every line parallel to $e_1$ meets $R^{-1}(N_v)$ in a null set, and by Tonelli's theorem
(Theorem~\ref{measure-theory:theorem-fubini}) $\lambda(N_v) = \lambda(R^{-1}(N_v)) = 0$. In particular the gradient $\nabla f = (D_{e_1}f, \ldots, D_{e_n}f)$ exists almost everywhere, and
$|\nabla f| \leq \sqrt nL$.

\emph{Step 2: $D_vf = \langle v, \nabla f\rangle$ almost everywhere.} Let $\varphi$ be $C^1$ with compact support. By translation invariance,
\[
\int\frac{f(x + tv) - f(x)}{t}\varphi(x)\,dx = -\int f(x)\frac{\varphi(x) - \varphi(x - tv)}{t}\,dx .
\]
Let $t = 1/j \to 0$. On the left the integrands converge almost everywhere to $D_vf\,\varphi$ and are bounded by $L|\varphi|$; on the right they converge to $f\,D_v\varphi = f\langle v, \nabla\varphi\rangle$
and are bounded by a constant times the indicator of a compact set, by the mean value inequality and the continuity of $f$. By dominated convergence,
$\int D_vf\,\varphi = -\int f\langle v, \nabla\varphi\rangle = -\sum_iv_i\int f\,\partial_i\varphi = \sum_iv_i\int D_{e_i}f\,\varphi$, the last equality being the first one for $v = e_i$. So
$\int(D_vf - \langle v, \nabla f\rangle)\varphi = 0$ for all $\varphi$, and Lemma~\ref{geometric-measure-theory:lemma-vanishing-integrals} applies.

\emph{Step 3: differentiability.} Let $(v_k)$ be a dense sequence in the unit sphere, and let $A$ be the set of points where $\nabla f$ and all $D_{v_k}f$ exist and
$D_{v_k}f = \langle v_k, \nabla f\rangle$; its complement is a countable union of null sets. Fix $x \in A$, and for a unit vector $v$ and $t \neq 0$ put
$Q(v, t) = (f(x + tv) - f(x))/t - \langle v, \nabla f(x)\rangle$. Then $|Q(v, t) - Q(v^{\prime}, t)| \leq (L + \sqrt nL)|v - v^{\prime}|$. Given $\varepsilon > 0$, by compactness of the
sphere there is $N$ such that every unit vector is within $\varepsilon/(2(1 + \sqrt n)L)$ of some $v_k$ with $k \leq N$, and there is $\delta > 0$ with $|Q(v_k, t)| < \varepsilon/2$ for
$0 < |t| < \delta$ and $k \leq N$. Then $|Q(v, t)| < \varepsilon$ for all unit vectors $v$ and $0 < |t| < \delta$, that is, $|f(x + h) - f(x) - \langle\nabla f(x), h\rangle| \leq \varepsilon|h|$ for
$|h| < \delta$. So $f$ is differentiable at $x$, with derivative $h \mapsto \langle\nabla f(x), h\rangle$.
\end{proof}

\begin{counterexample}\label{geometric-measure-theory:counterexample-rademacher}
\begin{enumerate}
\item Continuity is not enough: the Takagi function (Theorem~\ref{real-analysis:theorem-takagi}) is continuous and differentiable nowhere, so it is not Lipschitz on any interval.
\item The exceptional null set can be dense. Let $q_1, q_2, \ldots$ be an enumeration of $\QQ$ and $f(x) = \sum_j2^{-j}|x - q_j|$, a $1$-Lipschitz function on $\RR$. For $h > 0$,
\[\begin{aligned}
&\frac{f(q_k + h) - f(q_k)}{h} - \frac{f(q_k) - f(q_k - h)}{h} \\
&\qquad = \sum_j2^{-j}\,\frac{|q_k - q_j + h| + |q_k - q_j - h| - 2|q_k - q_j|}{h} \geq 2^{1-k},
\end{aligned}\]
since every term is $\geq 0$ by the triangle inequality and the term $j = k$ is $2^{1-k}$. If $f$ were differentiable at $q_k$ the left side would tend to $0$; so $f$ is differentiable at no
rational point.
\item Existence of all directional derivatives at a point does not imply differentiability there, even for Lipschitz functions: $f(x, y) = x^2y/(x^2 + y^2)$, $f(0, 0) = 0$, is Lipschitz,
has $D_vf(0) = f(v)$ for every unit vector $v$, and $v \mapsto f(v)$ is not linear. Step 3 of the proof is where the Lipschitz condition turns directional derivatives into a derivative.
\end{enumerate}
\end{counterexample}

\section{The area formula}\label{geometric-measure-theory:section-area-formula}

For a linear map $T \colon \RR^n \to \RR^n$, $\lambda^n(T(A)) = |\det T|\lambda^n(A)$. For a linear map $L \colon \RR^k \to \RR^m$ with $k \leq m$ the image of a cube is a
$k$-dimensional parallelepiped in $\RR^m$, and its $k$-dimensional volume is measured by $\mathcal{H}^k$. The factor replacing $|\det T|$ is the Jacobian of $L$.

\begin{definition}\label{geometric-measure-theory:definition-jacobian}
The \emph{Jacobian} of a linear map $L \colon \RR^k \to \RR^m$, $k \leq m$, is $JL = \sqrt{\det(L^TL)}$. For a map $f$ differentiable at $x$, $Jf(x) = J(Df(x))$.
\end{definition}

\begin{lemma}\label{geometric-measure-theory:lemma-jacobian}
Let $L \colon \RR^k \to \RR^m$ be linear, $k \leq m$.
\begin{enumerate}
\item $(JL)^2$ is the sum of the squares of the $k \times k$ minors of the matrix of $L$, and $JL = 0$ if and only if $L$ is not injective. For $k = m$, $JL = |\det L|$.
\item If $L$ is injective, $L = O \circ S$ with $S = (L^TL)^{1/2}$ positive definite symmetric on $\RR^k$, $\det S = JL$, and $O \colon \RR^k \to \RR^m$ an isometry onto $L(\RR^k)$.
\item $\mathcal{H}^k(L(A)) = JL\,\lambda^k(A)$ for every Lebesgue measurable $A \subseteq \RR^k$.
\end{enumerate}
\end{lemma}

\begin{proof}
(1) By Cauchy--Binet (Proposition~\ref{multilinear-algebra:proposition-minors}(2)) applied to $L^T$ and $L$, $\det(L^TL) = \sum_S\det(L_{S})^2$, where $L_S$ is the $k \times k$ submatrix with rows
$S$. By Proposition~\ref{multilinear-algebra:proposition-minors}(3), some minor is nonzero if and only if $L$ has rank $k$. For $k = m$, $\det(L^TL) = (\det L)^2$.

(2) $L^TL$ is symmetric with $\langle L^TLv, v\rangle = |Lv|^2 > 0$ for $v \neq 0$, so by the spectral theorem (Theorem~\ref{linear-algebra:theorem-spectral}) it has a positive definite square root $S$,
with $\det S = \sqrt{\det L^TL}$. Then $O = LS^{-1}$ satisfies $O^TO = S^{-1}L^TLS^{-1} = I$, so $|Ov| = |v|$, and $O(\RR^k) = L(\RR^k)$.

(3) If $L$ is not injective, $L(\RR^k)$ lies in a linear subspace $W$ of dimension $k - 1$; an isometry $\RR^{k-1} \to W$ and Theorem~\ref{geometric-measure-theory:theorem-hausdorff-lebesgue} give
$\mathcal{H}^{k-1}(W \cap B) < \infty$ for every ball $B$, so $\mathcal{H}^k(W) = 0$ by Proposition~\ref{geometric-measure-theory:proposition-hausdorff}(4), and both sides vanish. If $L$ is injective, $O$ is an
isometry from $\RR^k$ onto $L(\RR^k)$, which preserves $\mathcal{H}^k$, so $\mathcal{H}^k(L(A)) = \mathcal{H}^k(S(A)) = \lambda^k(S(A)) = |\det S|\lambda^k(A)$ by
Theorem~\ref{geometric-measure-theory:theorem-hausdorff-lebesgue} and Lemma~\ref{measure-theory:lemma-linear-change}; for Lebesgue measurable $A$ write $A = B \cup N$ with $B$ Borel and $N$ contained in a Borel null
set, which $S$ maps to a null set.
\end{proof}

So $JL$ is the $k$-dimensional volume of the parallelepiped spanned by $Le_1, \ldots, Le_k$, the image of the unit cube. For a curve, $J\gamma(t) = |\gamma^{\prime}(t)|$; for a parametrised
surface $x(u, v)$ in $\RR^3$, $Jx = |x_u \times x_v|$, since the three $2 \times 2$ minors of $Dx$ are the components of $x_u \times x_v$ up to sign.

\begin{theorem}[Area formula]\label{geometric-measure-theory:theorem-area-formula}
Let $U \subseteq \RR^k$ be open and $f \colon U \to \RR^m$ an injective $C^1$ map whose differential $Df(x)$ is injective at every point. Then for every Lebesgue measurable $A \subseteq U$ the set
$f(A)$ is $\mathcal{H}^k$-measurable, and
\[
\mathcal{H}^k(f(A)) = \int_AJf\,d\lambda^k .
\]
Consequently $\int_{f(U)}g\,d\mathcal{H}^k = \int_Ug(f(x))\,Jf(x)\,d\lambda^k(x)$ for every Borel function $g \geq 0$ on $\RR^m$.
\end{theorem}

\begin{proof}
For an injective linear $L \colon \RR^k \to \RR^m$ let $m(L) = \min_{|v| = 1}|Lv| > 0$. Fix $t > 1$.

\emph{Step 1: a decomposition of $U$.} Call a pair $(B, L)$ of an open ball $B \subseteq U$ with rational centre and radius and an injective linear map $L$ with rational matrix \emph{good} if, for all
$x, y \in B$,
\begin{equation}\label{geometric-measure-theory:equation-good-pair}
\begin{aligned}
&t^{-1}|L(x - y)| \leq |f(x) - f(y)| \leq t|L(x - y)| \qquad\text{and}\\
&t^{-k}JL \leq Jf(x) \leq t^kJL .
\end{aligned}
\end{equation}
Every $x_0 \in U$ lies in the ball of some good pair. Indeed, let $\mu_0 = m(Df(x_0)) > 0$ and $\eta = (1 - t^{-1})\mu_0/2$. Since $|m(L) - m(L^{\prime})| \leq \|L - L^{\prime}\|$, $J$ is
continuous, and $Jf(x_0) > 0$, there is $L$ with rational matrix such that $\|Df(x_0) - L\| < \min(\eta, \mu_0)/2$ and $t^{-k}JL < Jf(x_0) < t^kJL$; then $m(L) \geq \mu_0/2$, so
$\eta \leq (1 - t^{-1})m(L)$. By continuity of $Df$ and of $J$ there is a rational ball $B \ni x_0$ on which $\|Df - L\| \leq \eta$ and still $t^{-k}JL \leq Jf \leq t^kJL$. For $x, y \in B$, by the mean value inequality (Proposition~\ref{real-analysis:proposition-derivative-basic}(4))
applied to $f - L$ on the convex set $B$, $|f(x) - f(y) - L(x - y)| \leq \eta|x - y| \leq (1 - t^{-1})|L(x - y)|$, which gives
$t^{-1}|L(x - y)| \leq |f(x) - f(y)| \leq (2 - t^{-1})|L(x - y)| \leq t|L(x - y)|$. Enumerate the countably many good pairs as $(B_j, L_j)$, and let $E_j = B_j \setminus \bigcup_{i < j}B_i$; these are
disjoint Borel sets covering $U$.

\emph{Step 2: estimates on one piece.} Let $A \subseteq E_j$ be Lebesgue measurable. By \eqref{geometric-measure-theory:equation-good-pair} the map $L_j(x) \mapsto f(x)$, from $L_j(B_j)$ to $f(B_j)$, is
$t$-Lipschitz, and so is its inverse. By Proposition~\ref{geometric-measure-theory:proposition-hausdorff}(3) and Lemma~\ref{geometric-measure-theory:lemma-jacobian}(3),
\[\begin{aligned}
t^{-k}JL_j\,\lambda^k(A) &\leq \mathcal{H}^k(f(A)) \leq t^kJL_j\,\lambda^k(A), \\
t^{-k}JL_j\,\lambda^k(A) &\leq \int_AJf \leq t^kJL_j\,\lambda^k(A),
\end{aligned}\]
the second pair by \eqref{geometric-measure-theory:equation-good-pair}. Hence $t^{-2k}\int_AJf \leq \mathcal{H}^k(f(A)) \leq t^{2k}\int_AJf$; here $\mathcal{H}^k$ is the outer measure.

\emph{Step 3: measurability.} A Lebesgue measurable $A \subseteq U$ is $K \cup N$ with $K$ a countable union of compact sets and $N$ null (Theorem~\ref{measure-theory:theorem-lebesgue-measure}).
Then $f(K)$ is a countable union of compact sets, hence Borel, and $\mathcal{H}^k(f(N)) \leq \sum_j\mathcal{H}^k(f(N \cap E_j)) = 0$ by Step 2. So $f(A)$ differs from a Borel set by an
$\mathcal{H}^k$-null set, and it is $\mathcal{H}^k$-measurable (Theorem~\ref{measure-theory:theorem-caratheodory}).

\emph{Step 4.} For measurable $A \subseteq U$, the sets $f(A \cap E_j)$ are measurable and disjoint, since $f$ is injective. Summing Step 2 over $j$,
$t^{-2k}\int_AJf \leq \mathcal{H}^k(f(A)) \leq t^{2k}\int_AJf$, and $t > 1$ was arbitrary. For a Borel set $G \subseteq \RR^m$, applying this to $A = f^{-1}(G)$ gives the integral formula for
$g = \mathbf{1}_G$; it follows for simple functions by linearity, and for Borel $g \geq 0$ by monotone convergence (Theorem~\ref{measure-theory:theorem-monotone-convergence}).
\end{proof}

For $k = m$ the area formula is the change of variables formula (Theorem~\ref{measure-theory:theorem-change-of-variables}), with $\mathcal{H}^m = \lambda^m$. The proof is the same idea: on small pieces
$f$ is almost linear, and linear maps multiply measure by the Jacobian.

\begin{example}\label{geometric-measure-theory:example-area-formula}
\begin{enumerate}
\item \emph{Length of curves.} For an injective $C^1$ curve $\gamma \colon (a, b) \to \RR^m$ with $\gamma^{\prime} \neq 0$, $\mathcal{H}^1(\gamma((a, b))) = \int_a^b|\gamma^{\prime}(t)|\,dt$. So the length
of a curve, defined by parametrisations, is the one-dimensional Hausdorff measure of its image, which depends only on the set.
\item \emph{Area of graphs.} For $h \colon U \to \RR$ of class $C^1$, $U \subseteq \RR^k$ open, the graph map $f(x) = (x, h(x))$ is injective with injective differential, and
$Jf = \sqrt{\det(I + \nabla h\,\nabla h^T)} = \sqrt{1 + |\nabla h|^2}$, since $I + ww^T$ is the identity on $w^\perp$ and multiplies $w$ by $1 + |w|^2$. So the graph has area
$\int_U\sqrt{1 + |\nabla h|^2}\,d\lambda^k$.
\item \emph{The area of the sphere.} The upper hemisphere of $S^2$ is the graph of $h(x) = \sqrt{1 - |x|^2}$ over the open unit disc, with $|\nabla h|^2 = |x|^2/(1 - |x|^2)$ and
$Jf = (1 - |x|^2)^{-1/2}$. In polar coordinates (Example~\ref{measure-theory:example-gaussian}), $\mathcal{H}^2 = \int_0^{2\pi}\int_0^1 r(1 - r^2)^{-1/2}\,dr\,d\theta = 2\pi$. The equator has finite $\mathcal{H}^1$-measure by (1), so it is
$\mathcal{H}^2$-null (Proposition~\ref{geometric-measure-theory:proposition-hausdorff}(4)), and $\mathcal{H}^2(S^2) = 4\pi$.
\end{enumerate}
\end{example}

\begin{counterexample}\label{geometric-measure-theory:counterexample-area-injective}
Injectivity is needed: for $f(t) = (\cos t, \sin t)$ on $(0, 4\pi)$, $\int Jf = 4\pi$, while the image, the circle minus one point, has $\mathcal{H}^1$-measure $2\pi$. The integral counts the points of the
image with multiplicity: in general $\int_AJf = \int_{\RR^m}\#(A \cap f^{-1}(y))\,d\mathcal{H}^k(y)$ (Remark~\ref{geometric-measure-theory:remark-area-coarea}).
\end{counterexample}

The area formula also computes the area of surfaces of revolution, and gives the second theorem of Pappus. As in Proposition~\ref{measure-theory:proposition-revolution-volume}, the axis is the $z$-axis and
$(r, z)$ are the distance from the axis and the height.

\begin{proposition}[Pappus--Guldin theorem for areas]\label{geometric-measure-theory:proposition-pappus-area}
Let $\gamma = (\rho, \zeta) \colon [a, b] \to \RR^2$ be a $C^1$ curve with $\rho \geq 0$ on $[a, b]$ and $\rho > 0$ on $(a, b)$, injective and regular on $(a, b)$, and let
$\Sigma = \{(\rho(t)\cos\theta, \rho(t)\sin\theta, \zeta(t)) : t \in [a, b], \theta \in \RR\}$ be the surface obtained by rotating it around the $z$-axis. Then
\[
\mathcal{H}^2(\Sigma) = 2\pi\int_\gamma r\,ds = 2\pi\bar r\,L(\gamma),
\]
where $\int_\gamma r\,ds = \int_a^b\rho(t)|\gamma^{\prime}(t)|\,dt$ is a line integral of the first kind (Definition~\ref{real-analysis:definition-first-kind}) and $\bar r = L(\gamma)^{-1}\int_\gamma r\,ds$ is the distance of the
centroid of the curve from the axis.
\end{proposition}

\begin{proof}
Let $\Phi(t, \theta) = (\rho(t)\cos\theta, \rho(t)\sin\theta, \zeta(t))$ on $U = (a, b) \times (0, 2\pi)$. It is injective: $\rho(t) > 0$ and $\zeta(t)$ are recovered from $\Phi(t, \theta)$ as the distance from the axis and
the height, which determine $t$ because $\gamma$ is injective on $(a, b)$, and then $\theta \in (0, 2\pi)$ is determined. Its partial derivatives
$\Phi_t = (\rho^{\prime}\cos\theta, \rho^{\prime}\sin\theta, \zeta^{\prime})$ and $\Phi_\theta = (-\rho\sin\theta, \rho\cos\theta, 0)$ are orthogonal, with $|\Phi_t| = |\gamma^{\prime}(t)|$ and $|\Phi_\theta| = \rho(t)$. So
$D\Phi^TD\Phi = \operatorname{diag}(|\gamma^{\prime}|^2, \rho^2)$, $J\Phi = \rho|\gamma^{\prime}| > 0$ (Definition~\ref{geometric-measure-theory:definition-jacobian}), and $D\Phi$ is injective. By the area formula
(Theorem~\ref{geometric-measure-theory:theorem-area-formula}) and Tonelli's theorem, $\mathcal{H}^2(\Phi(U)) = \int_0^{2\pi}\int_a^b\rho(t)|\gamma^{\prime}(t)|\,dt\,d\theta = 2\pi\int_\gamma r\,ds$. The rest of $\Sigma$ is contained in the
meridian $\Phi([a, b] \times \{0\})$ and the two circles (or points) $\Phi(\{a\} \times \RR)$ and $\Phi(\{b\} \times \RR)$. These are images of compact intervals under $C^1$ maps, which are Lipschitz by the mean value
inequality, so they have finite $\mathcal{H}^1$-measure by Proposition~\ref{geometric-measure-theory:proposition-hausdorff}(3) and are $\mathcal{H}^2$-null by part (4) of the same proposition.
\end{proof}

\begin{example}\label{geometric-measure-theory:example-pappus-area}
\begin{enumerate}
\item \emph{The sphere.} The half circle $\gamma(t) = (R\sin t, -R\cos t)$, $t \in [0, \pi]$, has $\int_\gamma r\,ds = \int_0^\pi R\sin t \cdot R\,dt = 2R^2$, so the sphere of radius $R$ has area $4\pi R^2$; the centroid of the half
circle is at distance $2R/\pi$ from the axis, as in Example~\ref{real-analysis:example-first-kind}(1).
\item \emph{Archimedes' hat-box theorem.} The part of that sphere between the heights $z = c$ and $z = c + h$, with $-R \leq c < c + h \leq R$, comes from $t$ with $-R\cos t \in [c, c + h]$, and its area is
$2\pi\int R\sin t \cdot R\,dt = 2\pi R\bigl[-R\cos t\bigr] = 2\pi Rh$. So the area of a spherical zone depends only on its height, and equals the area of the corresponding zone of the circumscribed cylinder: projecting
the sphere horizontally onto the cylinder preserves areas. This is the principle behind Lambert's cylindrical equal-area map.
\item \emph{The torus.} The circle of radius $a$ centred at distance $c > a$ from the axis has centroid at its centre, so the torus has area $2\pi c \cdot 2\pi a = 4\pi^2ac$. (The parametrisation
$t \mapsto (c + a\cos t, a\sin t)$, $t \in [0, 2\pi]$, is injective on $(0, 2\pi)$, as the proposition requires.)
\end{enumerate}
\end{example}

\begin{counterexample}\label{geometric-measure-theory:counterexample-gabriel-horn}
\emph{Torricelli's trumpet}, also called Gabriel's horn: rotate the graph of $1/x$, $x \geq 1$, around the $x$-axis. By the disc method (Proposition~\ref{measure-theory:proposition-revolution-volume}(1), and monotone convergence
as the length tends to infinity) the solid has volume $\pi\int_1^\infty x^{-2}\,dx = \pi$. By Proposition~\ref{geometric-measure-theory:proposition-pappus-area}, applied to $\gamma(x) = (1/x, x)$ on $[1, T]$ and letting $T \to \infty$,
its surface has area
\[
2\pi\int_1^\infty\frac1x\sqrt{1 + x^{-4}}\,dx \geq 2\pi\int_1^\infty\frac{dx}{x} = \infty .
\]
So the trumpet can be filled with a finite amount of paint, but its inner surface cannot be painted. Torricelli's discovery (1643) of a solid of infinite extent and finite volume was a surprise to his
contemporaries. The paradox disappears once one notes that area and volume scale differently: a layer of paint of fixed thickness does not fit into the narrow end of the trumpet.
\end{counterexample}

\begin{remark}\label{geometric-measure-theory:remark-area-coarea}
\emph{The general area and coarea formulas.} By Rademacher's theorem a Lipschitz map $f \colon \RR^k \to \RR^m$ has a Jacobian almost everywhere, and the area formula holds in the generality
\[
\int_AJf\,d\lambda^k = \int_{\RR^m}\#\bigl(A \cap f^{-1}(y)\bigr)\,d\mathcal{H}^k(y) \qquad (k \leq m),
\]
for Lebesgue measurable $A$, where $\#$ is the number of points \cite[Section 3.3]{evans-gariepy}, \cite[Section 3.2]{federer-gmt}. For $k \geq m$ there is a dual statement, the \emph{coarea
formula}: with $Jf = \sqrt{\det(Df\,Df^T)}$,
\[
\int_AJf\,d\lambda^k = \int_{\RR^m}\mathcal{H}^{k-m}\bigl(A \cap f^{-1}(y)\bigr)\,d\lambda^m(y) \qquad (k \geq m),
\]
which slices $A$ by the level sets of $f$ instead of parametrising it \cite[Section 3.4]{evans-gariepy}. For $f(x) = |x|$ on $\RR^n$, $Jf = |\nabla f| = 1$ away from $0$, and the coarea formula
gives integration in spherical shells, $\int_{\RR^n}g\,d\lambda = \int_0^\infty\bigl(\int_{|x| = r}g\,d\mathcal{H}^{n-1}\bigr)dr$; for a $C^1$ function $u$ it gives $\int|\nabla u| =
\int_\RR\mathcal{H}^{n-1}(\{u = s\})\,ds$, the starting point of the isoperimetric and Sobolev inequalities. A set is \emph{$k$-rectifiable} if it is covered, up to an $\mathcal{H}^k$-null set,
by countably many Lipschitz images of subsets of $\RR^k$; these sets, which include all $k$-dimensional submanifolds, are the objects on which geometric measure theory does calculus, and
the area formula is the bridge between them and Lebesgue measure.
\end{remark}

\section{The coarea formula}\label{geometric-measure-theory:section-coarea}

The area formula measures a set by parametrising it. The coarea formula measures it by slicing. The model is Fubini's theorem: the volume of $A \subseteq \RR^n$ is the integral of the
$(n-1)$-dimensional volumes of the slices $A \cap \{x_n = t\}$. If the slices are the level sets $f^{-1}(t)$ of a function $f$ which are not parallel hyperplanes, the slab between the levels $t$ and
$t + dt$ has thickness about $dt/|\nabla f|$: where the levels are crowded, the slab is thin. So the volume of the slab is about $\mathcal{H}^{n-1}(\text{slice})\,dt/|\nabla f|$, and integrating
$|\nabla f|$ instead of $1$ compensates exactly:
\[
\int_A|\nabla f|\,d\lambda^n = \int_\RR\mathcal{H}^{n-1}\bigl(A \cap f^{-1}(t)\bigr)\,dt .
\]
We prove this for $C^1$ functions without critical points; the general statement is in Remark~\ref{geometric-measure-theory:remark-area-coarea}. Throughout, $n \geq 2$.

\begin{lemma}\label{geometric-measure-theory:lemma-normal-determinant}
Let $v_1, \ldots, v_{n-1} \in \RR^n$ be linearly independent, $V$ the $n \times (n - 1)$ matrix with these columns, $\nu$ a unit vector orthogonal to all $v_i$, and $w \in \RR^n$. Then
\[
|\det(v_1, \ldots, v_{n-1}, w)| = JV \cdot |\langle w, \nu\rangle| .
\]
\end{lemma}

In words: the volume of a parallelepiped is the $(n-1)$-dimensional volume of its base times its height. Recall that $JV = \sqrt{\det(V^TV)}$ is the $(n-1)$-dimensional volume of the
base (Lemma~\ref{geometric-measure-theory:lemma-jacobian}).

\begin{proof}
Write $w = u + \langle w, \nu\rangle\nu$ with $u$ in the span of the $v_i$. The determinant is linear in the last column and vanishes when the columns are dependent, so
$\det(V, w) = \langle w, \nu\rangle\det(V, \nu)$. For the matrix $M = (V, \nu)$, $M^TM$ is block diagonal with blocks $V^TV$ and $\langle\nu, \nu\rangle = 1$, because $V^T\nu = 0$. Hence
$(\det M)^2 = \det(M^TM) = \det(V^TV) = (JV)^2$.
\end{proof}

\begin{theorem}[Coarea formula]\label{geometric-measure-theory:theorem-coarea}
Let $U \subseteq \RR^n$ be open and $f \colon U \to \RR$ of class $C^1$ with $\nabla f(x) \neq 0$ for all $x \in U$. For every Borel set $A \subseteq U$ the function
$t \mapsto \mathcal{H}^{n-1}(A \cap f^{-1}(t))$ is Lebesgue measurable, and
\[
\int_A|\nabla f|\,d\lambda^n = \int_\RR\mathcal{H}^{n-1}\bigl(A \cap f^{-1}(t)\bigr)\,dt .
\]
More generally, for every Borel function $g \colon U \to [0, \infty]$,
\begin{equation}\label{geometric-measure-theory:equation-coarea}
\int_Ug\,|\nabla f|\,d\lambda^n = \int_\RR\Bigl(\int_{f^{-1}(t)}g\,d\mathcal{H}^{n-1}\Bigr)dt .
\end{equation}
\end{theorem}

\begin{proof}
\emph{Step 1: local coordinates adapted to $f$.} Let $x_0 \in U$. Some partial derivative $\partial_if(x_0)$ is nonzero. A permutation of the coordinates is an isometry of $\RR^n$, so it preserves
$\lambda^n$ (Lemma~\ref{measure-theory:lemma-linear-change}), $\mathcal{H}^{n-1}$ and $|\nabla f|$; so we may assume $\partial_nf(x_0) \neq 0$. The map
$\psi(x) = (x_1, \ldots, x_{n-1}, f(x))$ has $\det D\psi(x_0) = \partial_nf(x_0) \neq 0$. By the inverse function theorem (Theorem~\ref{real-analysis:theorem-inverse-function}) there are open sets
$V \ni x_0$ and $W$ such that $\psi \colon V \to W$ is a $C^1$-diffeomorphism. Let $\phi = \psi^{-1} \colon W \to V$, and write points of $W \subseteq \RR^{n-1} \times \RR$ as $(y, t)$. Since
$\psi(\phi(y, t)) = (y, t)$,
\[
f(\phi(y, t)) = t \qquad\text{for } (y, t) \in W .
\]
For fixed $t$ let $W_t = \{y : (y, t) \in W\}$, open in $\RR^{n-1}$, and $\phi_t = \phi(-, t) \colon W_t \to \RR^n$. It is injective and $C^1$, and $D\phi_t$ is injective, as its columns are the first
$n - 1$ columns of the invertible matrix $D\phi$. Its image is $V \cap f^{-1}(t)$: a point $x \in V$ has $f(x) = t$ if and only if $\psi(x) = (y, t)$ for some $y$.

\emph{Step 2: the Jacobian identity.} We claim $|\det D\phi(y, t)|\,|\nabla f(\phi(y, t))| = J\phi_t(y)$. Differentiating $f(\phi(y, t)) = t$ with the chain rule gives
$\langle\nabla f, \partial_{y_i}\phi\rangle = 0$ for $i < n$ and $\langle\nabla f, \partial_t\phi\rangle = 1$, all evaluated at $\phi(y, t)$ resp.\ $(y, t)$. So $\nu = \nabla f/|\nabla f|$ is a unit vector
orthogonal to the columns $\partial_{y_i}\phi$ of $D\phi_t$, and Lemma~\ref{geometric-measure-theory:lemma-normal-determinant} with $w = \partial_t\phi$ gives
\[
|\det D\phi| = J\phi_t\cdot|\langle\partial_t\phi, \nu\rangle| = \frac{J\phi_t}{|\nabla f|} .
\]

\emph{Step 3: the formula on $V$.} Let $A \subseteq V$ be Borel, and $h = \mathbf{1}_A \circ \phi$, a Borel function on $W$. By the change of variables formula
(Theorem~\ref{measure-theory:theorem-change-of-variables}) for $\phi \colon W \to V$ and Step 2,
\[\begin{aligned}
\int_A|\nabla f|\,d\lambda^n &= \int_Wh(y, t)\,|\nabla f(\phi(y, t))|\,|\det D\phi(y, t)|\,d\lambda^n(y, t) \\
&= \int_Wh(y, t)\,J\phi_t(y)\,d\lambda^n(y, t) .
\end{aligned}\]
The integrand, extended by $0$ outside $W$, is a nonnegative Borel function on $\RR^{n-1} \times \RR$. By Tonelli's theorem (Theorem~\ref{measure-theory:theorem-fubini} and
Remark~\ref{measure-theory:remark-lebesgue-product}) the function $t \mapsto \int_{W_t}h(y, t)J\phi_t(y)\,dy$ is measurable and the last integral equals
$\int_\RR\bigl(\int_{W_t}h(y, t)J\phi_t(y)\,dy\bigr)dt$. By the area formula (Theorem~\ref{geometric-measure-theory:theorem-area-formula}) for $\phi_t$ and the Borel set
$B_t = \{y \in W_t : \phi(y, t) \in A\}$, whose image is $A \cap f^{-1}(t)$ by Step 1,
\[
\int_{W_t}h(y, t)J\phi_t(y)\,dy = \int_{B_t}J\phi_t\,d\lambda^{n-1} = \mathcal{H}^{n-1}\bigl(A \cap f^{-1}(t)\bigr).
\]
This proves the theorem for Borel subsets of $V$.

\emph{Step 4: globalisation.} $U$ is a countable union of compact sets, each covered by finitely many of the open sets $V$ of Step 1; so $U = \bigcup_jV_j$ for countably many such sets. Let
$E_j = V_j \setminus \bigcup_{i < j}V_i$, disjoint Borel sets covering $U$. For a Borel set $A \subseteq U$, apply Step 3 to $A \cap E_j \subseteq V_j$ and sum over $j$. On the left,
$\sum_j\int_{A \cap E_j}|\nabla f| = \int_A|\nabla f|$ by monotone convergence. On the right, for every $t$ the sets $A \cap E_j \cap f^{-1}(t)$ are disjoint Borel sets with union
$A \cap f^{-1}(t)$, and $\mathcal{H}^{n-1}$ is a measure on Borel sets (Proposition~\ref{geometric-measure-theory:proposition-hausdorff}), so
$\mathcal{H}^{n-1}(A \cap f^{-1}(t)) = \sum_j\mathcal{H}^{n-1}(A \cap E_j \cap f^{-1}(t))$, a sum of measurable functions of $t$; integrate with monotone convergence
(Theorem~\ref{measure-theory:theorem-monotone-convergence}). Finally, \eqref{geometric-measure-theory:equation-coarea} is the first formula for $g = \mathbf{1}_A$; it follows for simple Borel functions by
linearity and for all Borel $g \geq 0$ by monotone convergence, applied on both sides.
\end{proof}

For $f(x) = x_n$ the coarea formula is Fubini's theorem. Its most important special case is integration in polar coordinates.

\begin{corollary}[Polar coordinates]\label{geometric-measure-theory:corollary-polar}
Let $S^{n-1} = \{x \in \RR^n : |x| = 1\}$ and $\omega_n = \lambda^n(B(0, 1))$. Then $\mathcal{H}^{n-1}(S^{n-1}) = n\omega_n$, and for every Borel function $g \colon \RR^n \to [0, \infty]$,
\[
\int_{\RR^n}g\,d\lambda^n = \int_0^\infty\Bigl(\int_{|x| = r}g\,d\mathcal{H}^{n-1}\Bigr)dr = \int_0^\infty r^{n-1}\Bigl(\int_{S^{n-1}}g(r\theta)\,d\mathcal{H}^{n-1}(\theta)\Bigr)dr .
\]
In particular, if $g(x) = h(|x|)$ is radial, $\int_{\RR^n}g\,d\lambda^n = n\omega_n\int_0^\infty h(r)r^{n-1}\,dr$.
\end{corollary}

\begin{proof}
Apply Theorem~\ref{geometric-measure-theory:theorem-coarea} to $f(x) = |x|$ on $U = \RR^n \setminus \{0\}$. Here $\nabla f(x) = x/|x|$, so $|\nabla f| = 1$, and the point $0$ is a null set. This gives the
first equality. For the second, $\theta \mapsto r\theta$ is a similarity with ratio $r$ from $S^{n-1}$ onto the sphere of radius $r$, so
$\mathcal{H}^{n-1}(\{|x| = r\} \cap G) = r^{n-1}\mathcal{H}^{n-1}(S^{n-1} \cap r^{-1}G)$ for Borel sets $G$ (Proposition~\ref{geometric-measure-theory:proposition-hausdorff}(3)), which is the second equality for
$g = \mathbf{1}_G$; pass to simple functions and to limits by monotone convergence. For $g = \mathbf{1}_{B(0, 1)}$ we get $\omega_n = \mathcal{H}^{n-1}(S^{n-1})\int_0^1r^{n-1}\,dr = \mathcal{H}^{n-1}(S^{n-1})/n$.
\end{proof}

\begin{example}\label{geometric-measure-theory:example-polar}
\begin{enumerate}
\item \emph{Spheres.} $\mathcal{H}^1(S^1) = 2\omega_2 = 2\pi$ and $\mathcal{H}^2(S^2) = 3\omega_3 = 3 \cdot \frac43\pi = 4\pi$, in agreement with Example~\ref{geometric-measure-theory:example-area-formula}. The
formula $\mathcal{H}^{n-1}(S^{n-1}) = n\omega_n$ says that the volume of the ball, $\omega_nR^n$, has derivative $n\omega_nR^{n-1}$ in $R$ equal to the area of the sphere of radius $R$: adding a thin shell of
thickness $dR$ adds volume (area)$\cdot dR$.
\item \emph{Integrability of powers.} For $a \in \RR$, $\int_{|x| < 1}|x|^{-a}\,dx = n\omega_n\int_0^1r^{n-1-a}\,dr$ is finite if and only if $a < n$, and $\int_{|x| > 1}|x|^{-a}\,dx = n\omega_n\int_1^\infty r^{n-1-a}\,dr$
is finite if and only if $a > n$. So $|x|^{-a}$ is never integrable on all of $\RR^n$, and the borderline exponent is the dimension, as for $\sum_{k \in \ZZ^n \setminus 0}|k|^{-a}$.
\item \emph{The volume of balls.} By Fubini's theorem and Example~\ref{measure-theory:example-gaussian}, $\int_{\RR^n}e^{-|x|^2}dx = (\int_\RR e^{-t^2}dt)^n = \pi^{n/2}$. By the corollary and the
substitution $u = r^2$ it also equals $n\omega_n\int_0^\infty e^{-r^2}r^{n-1}\,dr = \frac{n\omega_n}2\int_0^\infty e^{-u}u^{n/2 - 1}\,du = \frac{n\omega_n}2\Gamma(\frac n2)$. Hence
$\omega_n = \pi^{n/2}/\Gamma(\frac n2 + 1)$, which solves Exercise~\ref{measure-theory:exercise-ball-volume}: $\omega_1 = 2$, $\omega_2 = \pi$, $\omega_3 = \frac43\pi$, $\omega_4 = \frac{\pi^2}{2}$, and
$\omega_n \to 0$ very fast as $n \to \infty$.
\end{enumerate}
\end{example}

\begin{remark}\label{geometric-measure-theory:remark-coarea}
\begin{enumerate}
\item \emph{Critical points.} For a Lipschitz $f$ the coarea formula holds without excluding critical points \cite[Section 3.4]{evans-gariepy}. On the critical set $Z = \{\nabla f = 0\}$ the left side
vanishes, so the formula contains the statement that $\mathcal{H}^{n-1}(Z \cap f^{-1}(t)) = 0$ for almost every $t$. This is remarkable because the set of critical values $f(Z)$ can have positive
measure: Sard's theorem, which says it is null, needs $f$ of class $C^n$, and Whitney constructed a $C^1$ function on $\RR^2$ that is not constant on a connected arc of critical points
\cite{whitney-1935}; its critical values fill an interval, yet each level set meets $Z$ in an $\mathcal{H}^1$-null set.
\item \emph{Isoperimetry and Sobolev inequalities.} For $u \in C^1_c(\RR^n)$ the coarea formula reads $\int|\nabla u| = \int_\RR\mathcal{H}^{n-1}(\{u = t\})\,dt$ (in the general version), and
$\{u = t\}$ is, for almost every $t > 0$, the boundary of $\{u > t\}$. The \emph{isoperimetric inequality} $\mathcal{H}^{n-1}(\partial E) \geq n\omega_n^{1/n}\lambda^n(E)^{(n-1)/n}$, with equality for
balls, then gives the \emph{Sobolev inequality} $\|u\|_{n/(n-1)} \leq (n\omega_n^{1/n})^{-1}\|\nabla u\|_1$ with its best constant; conversely, the Sobolev inequality applied to smoothed indicator functions
gives back the isoperimetric inequality. This equivalence between a geometric and an analytic inequality is due to Federer and Fleming and to Maz'ya \cite{federer-fleming-1960}.
\end{enumerate}
\end{remark}

\section{Densities and rectifiability}\label{geometric-measure-theory:section-densities}

By the Lebesgue density theorem (Corollary~\ref{lebesgue-differentiation:corollary-density}) a measurable set $A \subseteq \RR^n$ fills almost all of every small ball centred at almost every one of its
points. For a set of dimension $s < n$ the analogous question is whether $\mathcal{H}^s(A \cap \overline{B}(x, r))$ is comparable to $\alpha(s)(2r)^s$, the $s$-dimensional ``volume'' of a ball of radius
$r$. For smooth curves and surfaces the answer is yes, with ratio tending to $1$; for the Cantor set the ratio oscillates forever. This difference turns out to be the difference between the sets on which
one can do calculus, the rectifiable ones, and the others. Throughout, $X$ is a metric space and $\overline{B}(x, r) = \{y : d(x, y) \leq r\}$.

\begin{definition}\label{geometric-measure-theory:definition-density}
For $A \subseteq X$, $s > 0$ and $x \in X$, the \emph{upper} and \emph{lower $s$-dimensional densities} of $A$ at $x$ are
\[
\Theta^{*s}(A, x) = \limsup_{r \to 0}\frac{\mathcal{H}^s(A \cap \overline{B}(x, r))}{\alpha(s)(2r)^s}, \qquad \Theta^s_*(A, x) = \liminf_{r \to 0}\frac{\mathcal{H}^s(A \cap \overline{B}(x, r))}{\alpha(s)(2r)^s}.
\]
If they are equal, the common value is the \emph{density} $\Theta^s(A, x)$.
\end{definition}

The normalisation is chosen so that flat pieces have density $1$. For $A \subseteq \RR^n$ Lebesgue measurable and $s = n$, $\mathcal{H}^n = \lambda^n$ and $\alpha(n)(2r)^n = \lambda^n(\overline{B}(x, r))$,
so by the density theorem $\Theta^n(A, x) = 1$ for almost every $x \in A$ and $\Theta^n(A, x) = 0$ for almost every $x \notin A$ (spheres are null, so closed and open balls give the same ratios).

\begin{example}\label{geometric-measure-theory:example-density-curve}
\emph{Curves have density $1$.} Let $\gamma \colon (a, b) \to \RR^m$ be an injective $C^1$ map with $\gamma^{\prime} \neq 0$ which is a homeomorphism onto its image $\Gamma$. Then
$\Theta^1(\Gamma, x) = 1$ for every $x \in \Gamma$. Let $x = \gamma(t_0)$, $c = |\gamma^{\prime}(t_0)| > 0$, and $0 < \varepsilon < 1$. Choose $\delta > 0$ with $[t_0 - \delta, t_0 + \delta] \subseteq (a, b)$
and $|\gamma^{\prime}(t) - \gamma^{\prime}(t_0)| \leq \varepsilon c$ there. By the mean value inequality, applied to $\gamma(t) - \gamma^{\prime}(t_0)(t - t_0)$,
\[
(1 - \varepsilon)c|t - t_0| \leq |\gamma(t) - x| \leq (1 + \varepsilon)c|t - t_0| \qquad (|t - t_0| \leq \delta),
\]
and $(1 - \varepsilon)c \leq |\gamma^{\prime}(t)| \leq (1 + \varepsilon)c$ there. Since $\gamma^{-1}$ is continuous at $x$, there is $r_0 \in (0, (1 - \varepsilon)c\delta]$ such that
$|\gamma(t) - x| \leq r_0$ implies $|t - t_0| \leq \delta$. For $r < r_0$ the set $I_r = \{t : |\gamma(t) - x| \leq r\}$ therefore satisfies
\[
\bigl[t_0 - \tfrac{r}{(1 + \varepsilon)c}, t_0 + \tfrac{r}{(1 + \varepsilon)c}\bigr] \subseteq I_r \subseteq \bigl[t_0 - \tfrac{r}{(1 - \varepsilon)c}, t_0 + \tfrac{r}{(1 - \varepsilon)c}\bigr].
\]
By the area formula (Theorem~\ref{geometric-measure-theory:theorem-area-formula}), $\mathcal{H}^1(\Gamma \cap \overline{B}(x, r)) = \int_{I_r}|\gamma^{\prime}|$ lies between $2r\frac{1 - \varepsilon}{1 + \varepsilon}$ and
$2r\frac{1 + \varepsilon}{1 - \varepsilon}$. As $\alpha(1) = 1$ and $\varepsilon$ was arbitrary, the ratio tends to $1$. The homeomorphism condition matters: for a curve shaped like the digit $6$, whose end
$\gamma(t)$, $t \to b$, runs into an interior point $x$ from one side, a small ball around $x$ contains the arc of the curve through $x$, of length about $2r$, and the end arc, of length about $r$, and $\Theta^1(\Gamma, x) = \frac32$.
\end{example}

\begin{example}\label{geometric-measure-theory:example-density-cantor}
\emph{The Cantor set has no density at $0$.} Let $C$ be the Cantor set, $s = \log 2/\log 3$, and $c = \mathcal{H}^s(C)$, so that $\alpha(s)/2 \leq c \leq \alpha(s)$
(Example~\ref{geometric-measure-theory:example-cantor-dimension}). For $k \geq 0$, $C \cap [0, 3^{-k}] = 3^{-k}C$ (the piece $C_w$ for $w = 1\cdots1$ of length $k$, in the notation of
Example~\ref{geometric-measure-theory:example-ifs}), so $\mathcal{H}^s(C \cap [0, 3^{-k}]) = 3^{-ks}c = 2^{-k}c$. Moreover $C$ does not meet the gap $(3^{-k-1}, 2 \cdot 3^{-k-1})$, so
$C \cap [0, 2 \cdot 3^{-k-1}]$ is $C \cap [0, 3^{-k-1}]$ plus the single point $2 \cdot 3^{-k-1}$. Hence, with $r = 3^{-k}$ and $r = 2 \cdot 3^{-k-1}$ respectively,
\[\begin{aligned}
\frac{\mathcal{H}^s(C \cap \overline{B}(0, 3^{-k}))}{\alpha(s)(2 \cdot 3^{-k})^s} &= \frac{2^{-k}c}{\alpha(s)2^s2^{-k}} = \frac{c}{\alpha(s)2^s}, \\
\frac{\mathcal{H}^s(C \cap \overline{B}(0, 2 \cdot 3^{-k-1}))}{\alpha(s)(4 \cdot 3^{-k-1})^s} &= \frac{2^{-k-1}c}{\alpha(s)4^s2^{-k-1}} = \frac{c}{\alpha(s)4^s}.
\end{aligned}\]
So $\Theta^{*s}(C, 0) \geq c/(\alpha(s)2^s)$ and $\Theta^s_*(C, 0) \leq c/(\alpha(s)4^s)$, which is smaller by the factor $2^s \approx 1.55$; as $c > 0$, the density does not exist. The same computation
works at the left endpoint $f_w(0)$ of every interval $I_w$: for $r < 3^{-|w|}$ the ball $\overline{B}(f_w(0), r)$ meets $C$ in $f_w(C \cap [0, 3^{|w|}r])$, because $C$ has a gap of length at least
$3^{-|w|}$ to the left of $I_w$ (Example~\ref{geometric-measure-theory:example-cantor-dimension}(2)), and $\mathcal{H}^s$ scales by $3^{-|w|s}$ under $f_w$. These endpoints are dense in $C$.
\end{example}

For a general set we cannot expect the density to exist, but the upper density is bounded below almost everywhere, by a short covering argument that uses nothing but the definition of $\mathcal{H}^s$.

\begin{theorem}\label{geometric-measure-theory:theorem-upper-density}
Let $A \subseteq X$ with $\mathcal{H}^s(A) < \infty$, $s > 0$. Then $\Theta^{*s}(A, x) \geq 2^{-s}$ for $\mathcal{H}^s$-almost every $x \in A$.
\end{theorem}

\begin{proof}
For $0 < t < 2^{-s}$ and $k \geq 1$ let
\[
A_{t, k} = \bigl\{x \in A : \mathcal{H}^s(A \cap \overline{B}(x, r)) \leq t\,\alpha(s)(2r)^s \text{ for all } 0 < r < 1/k\bigr\}.
\]
A point $x \in A$ with $\Theta^{*s}(A, x) < 2^{-s}$ lies in some $A_{t, k}$ with $t$ rational; so it suffices to show $\mathcal{H}^s(A_{t, k}) = 0$ for all $t$ and $k$, by countable subadditivity.

Fix $t, k$, let $0 < \delta < 1/k$ and $\varepsilon > 0$, and choose a cover $(C_i)$ of $A_{t, k}$ by sets of diameter at most $\delta$ with
$\sum_i\alpha(s)(\operatorname{diam}C_i)^s \leq \mathcal{H}^s_\delta(A_{t, k}) + \varepsilon$. We may discard the $C_i$ that do not meet $A_{t, k}$. For the others pick $x_i \in C_i \cap A_{t, k}$ and let
$d_i = \operatorname{diam}C_i$. If $d_i > 0$, then $C_i \subseteq \overline{B}(x_i, d_i)$ with $d_i < 1/k$, so by the definition of $A_{t, k}$
\[
\mathcal{H}^s(A_{t, k} \cap C_i) \leq \mathcal{H}^s(A \cap \overline{B}(x_i, d_i)) \leq t\,\alpha(s)(2d_i)^s = t2^s\,\alpha(s)d_i^s .
\]
If $d_i = 0$, then $C_i$ is a point and both sides are $0$. Summing, and using countable subadditivity of $\mathcal{H}^s$,
\[
\mathcal{H}^s(A_{t, k}) \leq \sum_i\mathcal{H}^s(A_{t, k} \cap C_i) \leq t2^s\bigl(\mathcal{H}^s_\delta(A_{t, k}) + \varepsilon\bigr) \leq t2^s\bigl(\mathcal{H}^s(A_{t, k}) + \varepsilon\bigr).
\]
Let $\varepsilon \to 0$. Since $\mathcal{H}^s(A_{t, k}) \leq \mathcal{H}^s(A) < \infty$ and $t2^s < 1$, this forces $\mathcal{H}^s(A_{t, k}) = 0$.
\end{proof}

The bound $2^{-s}$ comes from the comparison $C_i \subseteq \overline{B}(x_i, \operatorname{diam}C_i)$, which loses a factor $2$ in the radius. Inequalities in the other direction, bounding the
upper density from above, need a different idea: instead of starting from an arbitrary cover and enclosing its sets in balls, one starts from balls on which the measure is large and extracts disjoint
ones. This is done by the $5r$-covering lemma (Lemma~\ref{lebesgue-differentiation:lemma-five-r-covering}), which works in every metric space; recall that for a closed ball
$B = \overline{B}(x, r)$, given together with its centre and radius, $\widehat{B} = \overline{B}(x, 5r)$.

\begin{theorem}\label{geometric-measure-theory:theorem-density-lower-comparison}
Let $\mu$ be an outer measure on a metric space $X$ for which all Borel sets are $\mu$-measurable (for instance $\mu = \mathcal{H}^s$, by Lemma~\ref{geometric-measure-theory:lemma-metric-outer-measure}). Let
$A_1, A_2 \subseteq X$, $s > 0$ and $t \geq 0$, and suppose
\[
\Theta^{*s}(\mu, A_2, x) = \limsup_{r \to 0}\frac{\mu(A_2 \cap \overline{B}(x, r))}{\alpha(s)(2r)^s} \geq t \qquad \text{for all } x \in A_1 .
\]
Then $t\,\mathcal{H}^s(A_1) \leq \mu(A_2)$.
\end{theorem}

A Borel measure $\mu$ on $X$ defines such an outer measure, $\mu^*(E) = \inf\{\mu(B) : B \supseteq E \text{ Borel}\}$. The infimum is attained, by the intersection of Borel sets $B_j \supseteq E$ with
$\mu(B_j) \to \mu^*(E)$; and for $B$ Borel, $E \subseteq X$ and a Borel $E^{\prime} \supseteq E$ with $\mu(E^{\prime}) = \mu^*(E)$,
$\mu^*(E \cap B) + \mu^*(E \setminus B) \leq \mu(E^{\prime} \cap B) + \mu(E^{\prime} \setminus B) = \mu^*(E)$, so $B$ is $\mu^*$-measurable. So the theorem applies to Borel measures as well.
This is the form of the theorem in \cite{simon-gmt}, where $A_1 \subseteq A_2$ is assumed; the proof below does not use this inclusion. The sets $A_1$ and $A_2$ need not be measurable. The
important special case is $A_1 = A_2$; the freedom to choose $A_2$ is used to localise, with $A_2 = A \cap U$ for an open set $U \supseteq A_1$. Note the direction: Theorem~\ref{geometric-measure-theory:theorem-upper-density} and its generalisation Theorem~\ref{geometric-measure-theory:theorem-density-comparison} below bound $\mu$ from above where the density is small;
this theorem bounds $\mu$ from below where the density is large, and loses no factor $2^s$.

\begin{proof}
We may assume $t > 0$ and $\mu(A_2) < \infty$. Let $0 < t^{\prime} < t$ and $\delta > 0$, and let $\mathcal{F}$ be the family of closed balls $\overline{B}(x, r)$ with $x \in A_1$, $0 < r < \delta$ and
$\mu(A_2 \cap \overline{B}(x, r)) \geq t^{\prime}\alpha(s)(2r)^s$. By hypothesis every $x \in A_1$ is the centre of balls of $\mathcal{F}$ of arbitrarily small radius. Let $\mathcal{G}$ be as in
Lemma~\ref{lebesgue-differentiation:lemma-five-r-covering}.

The balls of $\mathcal{G}$ are disjoint Borel sets, hence $\mu$-measurable, and by the definition of measurability $\mu(A_2) \geq \sum_{B \in \mathcal{G}^{\prime}}\mu(A_2 \cap B)$ for every finite
$\mathcal{G}^{\prime} \subseteq \mathcal{G}$ (split $A_2$ successively along the balls). Each term is at least $t^{\prime}\alpha(s)(2r_B)^s > 0$, so only finitely many balls of $\mathcal{G}$ have $\mu(A_2 \cap B) > 1/m$, for each
$m$; hence $\mathcal{G}$ is countable, say $\mathcal{G} = \{B_1, B_2, \ldots\}$ with radii $r_i$, and
\[
\sum_it^{\prime}\alpha(s)(2r_i)^s \leq \sum_i\mu(A_2 \cap B_i) \leq \mu(A_2) < \infty .
\]

We claim that $A_1 \subseteq \bigcup_{i \leq N}B_i \cup \bigcup_{i > N}\widehat{B_i}$ for every $N$. Let $x \in A_1$ lie outside the closed set $B_1 \cup \cdots \cup B_N$. Some closed ball around $x$ misses this set, so
there is $B \in \mathcal{F}$ centred at $x$ and disjoint from $B_1, \ldots, B_N$. By the lemma $B$ meets some $B_i$ with $B \subseteq \widehat{B_i}$, and necessarily $i > N$; so $x \in \widehat{B_i}$.

The sets $B_i$ and $\widehat{B_i}$ have diameters at most $2r_i < 2\delta$ and $10r_i < 10\delta$. Hence
\[
\mathcal{H}^s_{10\delta}(A_1) \leq \sum_{i \leq N}\alpha(s)(2r_i)^s + \sum_{i > N}\alpha(s)(10r_i)^s \leq \frac{\mu(A_2)}{t^{\prime}} + 5^s\sum_{i > N}\alpha(s)(2r_i)^s .
\]
The last sum is the tail of a convergent series, so letting $N \to \infty$ gives $\mathcal{H}^s_{10\delta}(A_1) \leq \mu(A_2)/t^{\prime}$. Letting $\delta \to 0$ and then $t^{\prime} \to t$ proves the theorem.
\end{proof}

Together with Theorem~\ref{geometric-measure-theory:theorem-upper-density} this gives the basic density estimates for sets of finite Hausdorff measure, in every metric space.

\begin{theorem}[Upper density theorem]\label{geometric-measure-theory:theorem-density-bounds}
Let $A \subseteq X$ be $\mathcal{H}^s$-measurable with $\mathcal{H}^s(A) < \infty$, $s > 0$. Then $2^{-s} \leq \Theta^{*s}(A, x) \leq 1$ for $\mathcal{H}^s$-almost every $x \in A$, and $\Theta^{*s}(A, x) = 0$ for
$\mathcal{H}^s$-almost every $x \in X \setminus A$.
\end{theorem}

\begin{proof}
The lower bound is Theorem~\ref{geometric-measure-theory:theorem-upper-density}. Let $\nu(E) = \mathcal{H}^s(A \cap E)$ for Borel $E$; since $A$ and $E$ are $\mathcal{H}^s$-measurable, $\nu$ is a finite Borel measure, so
by Lemma~\ref{radon-measures:lemma-closed-approximation} every Borel set $E$ has open neighbourhoods $U$ with $\nu(U) < \nu(E) + \varepsilon$. For $x$ in an open set $U$, small balls around $x$ lie in $U$, so
$\Theta^{*s}(\mathcal{H}^s, A \cap U, x) = \Theta^{*s}(A, x)$.

\emph{Upper bound.} Let $t > 1$ and $B = \{x \in A : \Theta^{*s}(A, x) > t\}$. By Borel regularity (Remark~\ref{geometric-measure-theory:remark-hausdorff-conventions}(2)) there is a Borel set $E \supseteq B$ with
$\mathcal{H}^s(E) = \mathcal{H}^s(B)$, so $\nu(E) = \mathcal{H}^s(B)$, as $B \subseteq A \cap E \subseteq E$. For an open $U \supseteq E$ with $\nu(U) < \nu(E) + \varepsilon$,
Theorem~\ref{geometric-measure-theory:theorem-density-lower-comparison} with $\mu = \mathcal{H}^s$, $A_1 = B$, $A_2 = A \cap U$ gives $t\mathcal{H}^s(B) \leq \mathcal{H}^s(A \cap U) = \nu(U) < \mathcal{H}^s(B) + \varepsilon$. As
$\varepsilon$ is arbitrary and $\mathcal{H}^s(B) < \infty$, $(t - 1)\mathcal{H}^s(B) \leq 0$, so $\mathcal{H}^s(B) = 0$. Taking $t = 1 + 1/j$, $j \geq 1$, gives the upper bound.

\emph{Outside $A$.} Let $t > 0$ and $B = \{x \in X \setminus A : \Theta^{*s}(A, x) > t\}$. Choose a Borel set $A^{\prime} \supseteq A$ with $\mathcal{H}^s(A^{\prime}) = \mathcal{H}^s(A)$; since $A$ is measurable and of finite
measure, $\mathcal{H}^s(A^{\prime} \setminus A) = 0$, so $B \cap A^{\prime}$ is null. The Borel set $X \setminus A^{\prime}$ has $\nu(X \setminus A^{\prime}) = 0$, so it has open neighbourhoods $U$ with $\nu(U) < \varepsilon$, and
the theorem with $A_1 = B \setminus A^{\prime}$, $A_2 = A \cap U$ gives $t\mathcal{H}^s(B \setminus A^{\prime}) \leq \varepsilon$. Hence $B$ is null, and taking $t = 1/j$ gives the last statement.
\end{proof}

The name comes from \cite{simon-gmt}: the theorem is about upper densities only. In $\RR^n$ with $s = n$ it says less than the Lebesgue density theorem, which gives density exactly $1$ on $A$ and $0$
outside; but for non-integral $s$ the density need not exist (Example~\ref{geometric-measure-theory:example-density-cantor}), and the upper density is then the right object. For $A$ the Cantor set the theorem
says that $\Theta^{*s}(C, x) \in [2^{-s}, 1]$ at $\mathcal{H}^s$-almost every point of $C$.

So every set of finite $s$-dimensional measure is, at almost all of its points, at least somewhat dense and at most as dense as a flat piece. Whether the density exists, that is, whether the lower
density equals the upper one, is a deep geometric property, and it is related to the following class of sets.

\begin{definition}\label{geometric-measure-theory:definition-rectifiable}
Let $k \geq 1$ be an integer. A set $A \subseteq \RR^n$ is \emph{countably $k$-rectifiable} if there are Lipschitz maps $g_j \colon E_j \to \RR^n$, $E_j \subseteq \RR^k$, $j = 1, 2, \ldots$, such that
\[
\mathcal{H}^k\Bigl(A \setminus \bigcup_jg_j(E_j)\Bigr) = 0 .
\]
It is \emph{purely $k$-unrectifiable} if $\mathcal{H}^k(A \cap g(E)) = 0$ for every Lipschitz map $g \colon E \to \RR^n$ with $E \subseteq \RR^k$.
\end{definition}

So a rectifiable set is, up to a null set, a countable union of Lipschitz images of pieces of $\RR^k$, and on such pieces Rademacher's theorem and the area formula give tangent planes and a calculus. A
purely unrectifiable set meets each Lipschitz image of a piece of $\RR^k$ only in a null set.

\begin{proposition}\label{geometric-measure-theory:proposition-rectifiable}
\begin{enumerate}
\item Subsets and countable unions of countably $k$-rectifiable sets are countably $k$-rectifiable.
\item A countably $k$-rectifiable set $A$ has $\dim_{\mathcal{H}}A \leq k$, and $A \cap \bigcup_jg_j(E_j)$ has $\sigma$-finite $\mathcal{H}^k$-measure.
\item If $U \subseteq \RR^k$ is open and $g \colon U \to \RR^n$ is of class $C^1$, then $g(U)$ is countably $k$-rectifiable.
\item If $A$ is countably $k$-rectifiable and $B$ is purely $k$-unrectifiable, then $\mathcal{H}^k(A \cap B) = 0$.
\end{enumerate}
\end{proposition}

\begin{proof}
(1) is clear from the definition: a subset can use the same maps, and countably many countable families of maps form a countable family.

(2) Let $L_j$ be a Lipschitz constant of $g_j$. For $m \geq 1$, Proposition~\ref{geometric-measure-theory:proposition-hausdorff}(3) and Theorem~\ref{geometric-measure-theory:theorem-hausdorff-lebesgue} give
$\mathcal{H}^k(g_j(E_j \cap [-m, m]^k)) \leq L_j^k\mathcal{H}^k([-m, m]^k) = L_j^k(2m)^k < \infty$. So $\bigcup_jg_j(E_j)$ is a countable union of sets of finite $\mathcal{H}^k$-measure, each of which has
$\mathcal{H}^t$-measure $0$ for $t > k$ (Proposition~\ref{geometric-measure-theory:proposition-hausdorff}(4)); the remaining part of $A$ is $\mathcal{H}^k$-null, hence also $\mathcal{H}^t$-null for $t > k$.

(3) $U$ is the union of the countably many closed balls $\overline{B} \subseteq U$ with rational centres and radii. On such a ball, which is compact and convex, $\|Dg\|$ is bounded, say by $L$, and the mean
value inequality (Proposition~\ref{real-analysis:proposition-derivative-basic}(4)) shows that $g$ is $L$-Lipschitz on $\overline{B}$.

(4) $A \cap B$ is contained in the union of the $\mathcal{H}^k$-null set $A \setminus \bigcup_jg_j(E_j)$ and the sets $B \cap g_j(E_j)$, which are $\mathcal{H}^k$-null by assumption.
\end{proof}

By (3), images of $C^1$ maps, such as the curves of Example~\ref{geometric-measure-theory:example-density-curve}, graphs of $C^1$ or Lipschitz functions, and $k$-dimensional submanifolds of $\RR^n$ (which are
covered by countably many images of charts), are countably $k$-rectifiable. At the other extreme, the four-corner Cantor set $K_4$ of Example~\ref{geometric-measure-theory:example-similarity-dimension}(2), which has
$0 < \mathcal{H}^1(K_4) < \infty$, is purely $1$-unrectifiable \cite{falconer-gofs}, \cite{mattila-1995}. The two classes are distinguished by densities and by projections, and these are among the
deepest results of the subject.

\begin{theorem}\label{geometric-measure-theory:theorem-rectifiability-criteria}
Let $A \subseteq \RR^n$ be $\mathcal{H}^k$-measurable with $0 < \mathcal{H}^k(A) < \infty$, $k$ an integer with $1 \leq k < n$.
\begin{enumerate}
\item (\emph{Densities}) $A$ is countably $k$-rectifiable if and only if $\Theta^k(A, x) = 1$ for $\mathcal{H}^k$-almost every $x \in A$, and if and only if $\Theta^k(A, x)$ exists and is positive for
$\mathcal{H}^k$-almost every $x \in A$.
\item (\emph{Projections, Besicovitch--Federer}) If $A$ is countably $k$-rectifiable, then its orthogonal projections to almost all $k$-dimensional linear subspaces have positive $\mathcal{H}^k$-measure. If
$A$ is purely $k$-unrectifiable, then its projections to almost all $k$-dimensional linear subspaces are $\mathcal{H}^k$-null.
\end{enumerate}
Here ``almost all'' refers to the rotation-invariant measure on the set of $k$-dimensional subspaces; for $k = 1$, $n = 2$ it means almost all angles. See \cite{mattila-1995}, \cite{federer-gmt}; the last
statement of (1) is Preiss's theorem \cite{preiss-1987}.
\end{theorem}

\begin{remark}\label{geometric-measure-theory:remark-densities}
\begin{enumerate}
\item \emph{Non-integral dimensions.} If $0 < \mathcal{H}^s(A) < \infty$ and the density $\Theta^s(A, x)$ exists and is positive at almost every point of $A$, then $s$ is an integer (Marstrand, see
\cite{mattila-1995}). So sets of non-integral dimension, such as the Cantor set (Example~\ref{geometric-measure-theory:example-density-cantor}) and the self-similar sets of
Section~\ref{geometric-measure-theory:section-self-similar}, are always irregular: their densities oscillate.
\item \emph{The four-corner set again.} Example~\ref{geometric-measure-theory:example-four-corner-projections} computed two projections of $K_4$: in the direction of the axes it has length $0$, as
Theorem~\ref{geometric-measure-theory:theorem-rectifiability-criteria}(2) predicts for almost all directions, while in the direction of $(2, 1)$ it is an interval. So the exceptional set of directions in the
Besicovitch--Federer theorem is not empty. For a segment the situation is reversed: its projection has positive length in every direction except the perpendicular one.
\item \emph{Besicovitch and Kakeya sets.} A \emph{Besicovitch set} is a compact subset of $\RR^n$ that contains a segment of length $1$ in every direction. Besicovitch showed that such sets exist in
$\RR^2$ with Lebesgue measure $0$; correspondingly, a needle of length $1$ can be turned around completely inside a region of arbitrarily small area (\emph{Kakeya's needle problem}). Every Besicovitch
set in $\RR^2$ has Hausdorff dimension $2$ (Davies), see \cite{falconer-gofs}. Whether Besicovitch sets in $\RR^n$ must have Hausdorff dimension $n$ is the \emph{Kakeya conjecture}, a central problem
linking geometric measure theory with harmonic analysis and analytic number theory.
\end{enumerate}
\end{remark}

\subsection{Densities of measures}\label{geometric-measure-theory:subsection-measure-densities}

Nothing in the definition of densities is special to the Hausdorff measure of a set: one can ask how much mass any measure puts in small balls. This point of view, in which the set $A$ is replaced
by the measure $\mathcal{H}^s$ restricted to $A$, is the one used in modern geometric measure theory, and it simplifies the statements. Throughout, $\mu$ is a Borel measure on the metric space $X$ that is
finite on bounded Borel sets.

\begin{definition}\label{geometric-measure-theory:definition-measure-density}
For $s > 0$ and $x \in X$, the \emph{upper} and \emph{lower $s$-dimensional densities} of $\mu$ at $x$ are
\[
\Theta^{*s}(\mu, x) = \limsup_{r \to 0}\frac{\mu(\overline{B}(x, r))}{\alpha(s)(2r)^s}, \qquad \Theta^s_*(\mu, x) = \liminf_{r \to 0}\frac{\mu(\overline{B}(x, r))}{\alpha(s)(2r)^s},
\]
and the \emph{density} $\Theta^s(\mu, x)$ is their common value when they agree.
\end{definition}

\begin{example}\label{geometric-measure-theory:example-measure-density}
\begin{enumerate}
\item If $A \subseteq X$ is a Borel set with $\mathcal{H}^s(A) < \infty$, then $\mu = \mathcal{H}^s\llcorner A$, $\mu(B) = \mathcal{H}^s(A \cap B)$, is such a measure (Borel sets are
$\mathcal{H}^s$-measurable by Lemma~\ref{geometric-measure-theory:lemma-metric-outer-measure}), and $\Theta^{*s}(\mu, x) = \Theta^{*s}(A, x)$: the densities of Definition~\ref{geometric-measure-theory:definition-density} are those
of this measure.
\item For $f \in L^1_{\mathrm{loc}}(\RR^n)$ with $f \geq 0$ and $\mu = f\lambda$, $\Theta^n(\mu, x) = f(x)$ for almost every $x$, since $\alpha(n)(2r)^n = \lambda(\overline{B}(x, r))$ and the averages of $f$ over balls
converge to $f(x)$ almost everywhere (Theorem~\ref{lebesgue-differentiation:theorem-lebesgue-differentiation}; closed and open balls differ by a null sphere). For $s < n$, $\Theta^s(\mu, x) = 0$ wherever $f$ is
bounded near $x$.
\item For a point mass $\mu = \delta_a$, $\Theta^{*s}(\mu, a) = \infty$ for every $s > 0$, and $\Theta^s(\mu, x) = 0$ for $x \neq a$.
\item For the Cantor measure $\mu$ of Example~\ref{geometric-measure-theory:example-cantor-dimension}(3) and $s = \log 2/\log 3$, part (4) of that example gives $\mu(\overline{B}(x, r)) \leq 2(2r)^s$ for $2r < 1$, so
$\Theta^{*s}(\mu, x) \leq 2/\alpha(s)$ at every point. At $0$ the density does not exist: $\mu([0, 3^{-k}]) = 2^{-k}$, and $\mu([0, 2 \cdot 3^{-k-1}]) = 2^{-k-1}$ because $\mu$ gives no
mass to the gap $(3^{-k-1}, 2 \cdot 3^{-k-1})$ and none to points (the Cantor function is continuous), so the computation of Example~\ref{geometric-measure-theory:example-density-cantor} applies with $c = 1$.
\end{enumerate}
\end{example}

The basic theorem compares $\mu$ with $\mathcal{H}^s$ on the set where the density of $\mu$ is bounded.

\begin{theorem}\label{geometric-measure-theory:theorem-density-comparison}
Let $\mu$ be a Borel measure on $X$, finite on bounded Borel sets, let $A \subseteq X$ be a Borel set, $s > 0$ and $0 < t < \infty$. If $\Theta^{*s}(\mu, x) \leq t$ for every $x \in A$, then
\[
\mu(A) \leq 2^st\,\mathcal{H}^s(A).
\]
\end{theorem}

\begin{proof}
\emph{Measurability.} For fixed $r > 0$ the function $x \mapsto \mu(\overline{B}(x, r))$ is upper semicontinuous: if $x_j \to x$, then $\overline{B}(x_j, r) \subseteq \overline{B}(x, r + \eta)$ for large $j$,
so $\limsup_j\mu(\overline{B}(x_j, r)) \leq \mu(\overline{B}(x, r + \eta))$ for every $\eta > 0$, and $\mu(\overline{B}(x, r + \eta)) \to \mu(\overline{B}(x, r))$ as $\eta \to 0$ by continuity from above, since
$\overline{B}(x, r) = \bigcap_{\eta > 0}\overline{B}(x, r + \eta)$ and these balls have finite measure. So the sets $\{x : \mu(\overline{B}(x, r)) < c\}$ are open and the sets
$\{x : \mu(\overline{B}(x, r)) \leq c\}$ are Borel. The same continuity shows that $r \mapsto \mu(\overline{B}(x, r))$ is right continuous.

\emph{Covering.} Let $t^{\prime} > t$ and, for $k \geq 1$,
\[
A_k = \bigl\{x \in A : \mu(\overline{B}(x, r)) \leq t^{\prime}\alpha(s)(2r)^s \text{ for all } 0 < r < 1/k\bigr\}.
\]
By right continuity it suffices to require the inequality for rational $r$, so $A_k$ is Borel. The sets $A_k$ increase, and their union is $A$, because $\Theta^{*s}(\mu, x) \leq t < t^{\prime}$. Fix $k$,
$0 < \delta < 1/k$ and $\varepsilon > 0$, and assume $\mathcal{H}^s(A_k) < \infty$ (otherwise there is nothing to prove). Choose a cover of $A_k$ by closed sets $C_i$ of diameter $d_i \leq \delta$ with
$\sum_i\alpha(s)d_i^s \leq \mathcal{H}^s_\delta(A_k) + \varepsilon$, discard those not meeting $A_k$, and pick $x_i \in C_i \cap A_k$. If $d_i > 0$, then $C_i \subseteq \overline{B}(x_i, d_i)$ with $d_i < 1/k$, so
$\mu(C_i) \leq t^{\prime}\alpha(s)(2d_i)^s = 2^st^{\prime}\alpha(s)d_i^s$. If $d_i = 0$, then $C_i = \{x_i\}$ and $\mu(C_i) \leq \mu(\overline{B}(x_i, r)) \leq t^{\prime}\alpha(s)(2r)^s$ for all small $r$, so
$\mu(C_i) = 0$. Summing,
\[
\mu(A_k) \leq \sum_i\mu(C_i) \leq 2^st^{\prime}\bigl(\mathcal{H}^s_\delta(A_k) + \varepsilon\bigr) \leq 2^st^{\prime}\bigl(\mathcal{H}^s(A) + \varepsilon\bigr).
\]
Letting $\varepsilon \to 0$, then $k \to \infty$ (continuity from below), then $t^{\prime} \to t$ gives the claim.
\end{proof}

For $\mu = \mathcal{H}^s\llcorner A$ the argument is the proof of Theorem~\ref{geometric-measure-theory:theorem-upper-density}: on the set where $\Theta^{*s}(A, \cdot) \leq t$ with $t < 2^{-s}$ the theorem
gives $\mathcal{H}^s \leq 2^st\mathcal{H}^s$, which forces measure zero. The hypothesis $t > 0$ cannot be dropped: Lebesgue measure on $\RR$ has $\Theta^{1/2}(\lambda, x) = 0$ everywhere, while
$\lambda(\RR) > 0$ and $\mathcal{H}^{1/2}(\RR) = \infty$.

\begin{corollary}\label{geometric-measure-theory:corollary-density-consequences}
Let $\mu$ be as in the theorem and $s > 0$.
\begin{enumerate}
\item (\emph{Absolute continuity}) If $\Theta^{*s}(\mu, x) < \infty$ for every $x$ in a Borel set $A$, then every Borel set $E \subseteq A$ with $\mathcal{H}^s(E) = 0$ has $\mu(E) = 0$.
\item (\emph{Mass distribution principle}) If $A$ is a Borel set and there are $c, r_0 > 0$ with $\mu(\overline{B}(x, r)) \leq cr^s$ for all $x \in A$ and $0 < r < r_0$, then
$\mathcal{H}^s(A) \geq \alpha(s)\mu(A)/c$. In particular, if $\mu(A) > 0$ then $\dim_{\mathcal{H}}A \geq s$.
\end{enumerate}
\end{corollary}

\begin{proof}
(1) For integers $m, k \geq 1$ let $E_{m,k}$ be the set of $x \in E$ with $\mu(\overline{B}(x, r)) \leq m\alpha(s)(2r)^s$ for $0 < r < 1/k$, a Borel set as in the proof above, on which
$\Theta^{*s}(\mu, \cdot) \leq m$. By the theorem $\mu(E_{m,k}) \leq 2^sm\mathcal{H}^s(E_{m,k}) = 0$, and $E = \bigcup_{m,k}E_{m,k}$ because the upper density is finite on $E$.
(2) On $A$, $\Theta^{*s}(\mu, x) \leq c/(2^s\alpha(s))$, so the theorem gives $\mu(A) \leq (c/\alpha(s))\mathcal{H}^s(A)$. If $\mu(A) > 0$, then $\mathcal{H}^s(A) > 0$, so $\dim_{\mathcal{H}}A \geq s$.
\end{proof}

For the Cantor measure, (2) with $c = 2^{1+s}$ gives $\mathcal{H}^s(C) \geq \alpha(s)2^{-1-s}$, a little weaker than the bound $\alpha(s)/2$ of Example~\ref{geometric-measure-theory:example-cantor-dimension}, which
estimated $\mu$ on arbitrary small sets rather than on balls. Statement (1) says that a measure with finite upper $s$-density cannot charge $\mathcal{H}^s$-null sets; for instance a measure with finite
upper $1$-density everywhere on $\RR^2$ gives no mass to the Cantor set $C \times \{0\}$, of dimension $\log 2/\log 3 < 1$.

\begin{remark}\label{geometric-measure-theory:remark-measure-densities}
\begin{enumerate}
\item \emph{The reverse inequality.} Theorem~\ref{geometric-measure-theory:theorem-density-lower-comparison} is the reverse of the theorem above: for $\mu$ as in the theorem, a Borel set $A$ and an open
$U \supseteq A$ with $t_1 \leq \Theta^{*s}(\mu, x) \leq t_2$ on $A$, they give $t_1\mathcal{H}^s(A) \leq \mu(U)$ (as $\mu(U \cap \overline{B}(x, r)) = \mu(\overline{B}(x, r))$ for $x \in A$ and small $r$)
and $\mu(A) \leq 2^st_2\mathcal{H}^s(A)$; taking the infimum over $U$, which is $\mu(A)$ by outer regularity (Lemma~\ref{radon-measures:lemma-closed-approximation}, applied on bounded pieces),
$t_1\mathcal{H}^s(A) \leq \mu(A) \leq 2^st_2\mathcal{H}^s(A)$. The factor $2^s$ on the right is lost because the proof encloses arbitrary sets of a cover in balls; the left inequality is proved with
the $5r$-covering lemma, which works with balls from the start. For $\mu = \mathcal{H}^s\llcorner A$ the two inequalities are the two halves of the upper density theorem
(Theorem~\ref{geometric-measure-theory:theorem-density-bounds}).
\item \emph{Differentiation of one measure by another.} For two Radon measures $\mu$, $\nu$ on $\RR^n$, the limit $\lim_{r \to 0}\nu(\overline{B}(x, r))/\mu(\overline{B}(x, r))$ exists for $\mu$-almost every $x$
and is the Radon--Nikodym derivative of the absolutely continuous part of $\nu$ with respect to $\mu$ (Theorem~\ref{radon-measures:theorem-besicovitch-differentiation}). For $\mu = \lambda$ this is
Theorem~\ref{radon-measures:theorem-differentiation-measures}. For a general $\mu$ the Vitali covering lemma is useless, since $\mu(B(x, 3r))$ need not be comparable to $\mu(B(x, r))$, and it is
replaced by the Besicovitch covering theorem (Theorem~\ref{radon-measures:theorem-besicovitch}), which uses the geometry of $\RR^n$ and produces covers that split into $N(n)$ disjoint
families. In a general metric space the Vitali argument survives
when $\mu$ is doubling, $\mu(\overline{B}(x, 2r)) \leq C\mu(\overline{B}(x, r))$; see Subsection~\ref{radon-measures:subsection-doubling}.
\item \emph{Rectifiable measures.} The rectifiability criteria have their natural form for measures. A Radon measure $\mu$ on $\RR^n$ is \emph{$k$-rectifiable} if $\mu = \theta\mathcal{H}^k\llcorner E$
with $E$ countably $k$-rectifiable and $\theta$ a positive Borel function. \emph{Preiss's theorem} says that $\mu$ is $k$-rectifiable if and only if $0 < \Theta^k(\mu, x) < \infty$ for $\mu$-almost every
$x$ \cite{preiss-1987}, and \emph{Marstrand's theorem} says that if $0 < \Theta^s(\mu, x) < \infty$ on a set of positive $\mu$-measure, then $s$ is an integer \cite{mattila-1995}. Preiss's proof
studies the \emph{tangent measures} of $\mu$ at $x$, the weak limits (Section~\ref{radon-measures:section-weak-convergence}) of the rescaled measures $A \mapsto \mu(x + rA)/r^k$ as $r \to 0$; existence
of the density forces them to be flat.
\item \emph{Frostman's lemma.} The mass distribution principle has a converse: if $A \subseteq \RR^n$ is a Borel set with $\mathcal{H}^s(A) > 0$, there is a nonzero finite Borel measure $\mu$ supported on a
compact subset of $A$ with $\mu(\overline{B}(x, r)) \leq r^s$ for all $x$ and $r$ \cite{mattila-1995}. So $\dim_{\mathcal{H}}A$ is the supremum of the $s$ for which $A$ carries a nonzero measure with
$\mu(\overline{B}(x, r)) \leq r^s$: Hausdorff dimension can be computed with measures instead of covers.
\end{enumerate}
\end{remark}

\begin{remark}\label{geometric-measure-theory:remark-doubling-densities}
The densities of this section are related to the differentiation with respect to doubling measures (Section~\ref{radon-measures:subsection-doubling}) as follows.
\begin{enumerate}
\item \emph{Densities as derivatives.} The density $\Theta^s(\mu, x)$ is the derivative of $\mu$ with respect to the function $r \mapsto \alpha(s)(2r)^s$ in place of a measure $\nu$. Since this function
is not the measure of the balls of a fixed measure, there is no Lebesgue decomposition to recover, and one only gets the inequalities of Theorems~\ref{geometric-measure-theory:theorem-density-lower-comparison}
and~\ref{geometric-measure-theory:theorem-density-comparison}. The proofs are parallel: the singular part in Theorem~\ref{radon-measures:theorem-doubling-differentiation} is handled exactly as in Theorem~\ref{geometric-measure-theory:theorem-density-lower-comparison}, with
$\nu(B)$ in place of $\alpha(s)(2r)^s$ for $B = \overline{B}(x, r)$, and with \eqref{radon-measures:equation-doubling-five} in place of $\alpha(s)(10r)^s = 5^s\alpha(s)(2r)^s$.
\item \emph{Ahlfors regularity.} A measure $\nu$ with $c^{-1}r^s \leq \nu(\overline{B}(x, r)) \leq cr^s$ for all $x \in X$ and $0 < r \leq \operatorname{diam}X$ (so $c \geq 1$) is doubling with
$C = 2^sc^2$: if $2r \leq \operatorname{diam}X$, then $\nu(\overline{B}(x, 2r)) \leq c2^sr^s \leq 2^sc^2\nu(\overline{B}(x, r))$; if $r \leq \operatorname{diam}X < 2r$, then
$\nu(\overline{B}(x, 2r)) = \nu(\overline{B}(x, \operatorname{diam}X)) \leq c(2r)^s \leq 2^sc^2\nu(\overline{B}(x, r))$; and if $r > \operatorname{diam}X$, both balls are $X$. For such a measure, $\nu$ and
$\mathcal{H}^s$ are comparable on Borel sets $A$: $\nu(A) \leq (c/\alpha(s))\mathcal{H}^s(A)$ by Theorem~\ref{geometric-measure-theory:theorem-density-comparison}, and
$\mathcal{H}^s(A) \leq 2^sc\,\alpha(s)\nu(U)$ for every open $U \supseteq A$ by Theorem~\ref{geometric-measure-theory:theorem-density-lower-comparison} with $A_1 = A$, $A_2 = U$, which gives
$\mathcal{H}^s(A) \leq 2^sc\,\alpha(s)\nu(A)$ by outer regularity (Lemma~\ref{radon-measures:lemma-closed-approximation}, applied on bounded pieces). The Cantor measure on the Cantor set $C$ is of this
kind with $s = \log 2/\log 3$: Example~\ref{geometric-measure-theory:example-measure-density}(4) gives the upper bound, and for $x \in C$ and $3^{-k} \leq r < 3^{1-k}$ the ball $\overline{B}(x, r)$ contains the
interval of the $k$-th stage of the construction that contains $x$, of measure $2^{-k} = (3^{-k})^s \geq 3^{-s}r^s$.
\end{enumerate}
\end{remark}

\section{Exercises}\label{geometric-measure-theory:section-exercises}

\begin{exercise}\label{geometric-measure-theory:exercise-middle-half-cantor}
Let $K \subseteq [0, 1]$ be the set obtained by removing the open middle half $(\frac14, \frac34)$ of $[0, 1]$, then the open middle halves of the two remaining intervals, and so on; equivalently, the
numbers $\sum_{k \geq 1}a_k4^{-k}$ with all $a_k \in \{0, 3\}$. Show that $\dim_{\mathcal{H}}K = \frac12$.
\end{exercise}

\begin{solution}
Follow Example~\ref{geometric-measure-theory:example-cantor-dimension} with $3$ replaced by $4$ and the digits $\{0, 2\}$ by $\{0, 3\}$; now $s = \log 2/\log 4 = \frac12$, so $4^s = 2$. The $2^k$ intervals of
generation $k$ have length $4^{-k}$, which gives $\mathcal{H}^{1/2}(K) \leq \alpha(\frac12)2^k4^{-k/2} = \alpha(\frac12)$. Two of them first differing at position $j \leq k$ have left ends differing by at least
$3 \cdot 4^{-j} - \sum_{j < i \leq k}3 \cdot 4^{-i} \geq 2 \cdot 4^{-j}$, so they are at distance at least $4^{-k}$. The function $F(\sum_ka_k4^{-k}) = \sum_k(a_k/3)2^{-k}$ on $K$, extended to
$[0, 1]$ as in Counterexample~\ref{lebesgue-differentiation:counterexample-cantor-function}, is continuous and nondecreasing, and its measure $\mu_F$ (Lemma~\ref{lebesgue-differentiation:lemma-stieltjes}) gives mass
$2^{-k}$ to each interval of generation $k$ and mass $0$ to $[0, 1] \setminus K$. A closed set of diameter $d$ with $4^{-k-1} \leq d < 4^{-k}$ meets at most one interval of generation $k$, so its
$\mu_F$-measure is at most $2^{-k} = 2 \cdot 4^{-(k+1)/2} \leq 2d^{1/2}$, and as in the example $\mathcal{H}^{1/2}(K) \geq \alpha(\frac12)/2 > 0$.
\end{solution}

\begin{exercise}\label{geometric-measure-theory:exercise-lipschitz-graph}
Let $h \colon \RR^n \to \RR$ be Lipschitz with constant $L$ and $U \subseteq \RR^n$ a bounded open set. Show that the graph $\Gamma = \{(x, h(x)) : x \in U\}$ satisfies
$\lambda^n(U) \leq \mathcal{H}^n(\Gamma) \leq (1 + L^2)^{n/2}\lambda^n(U)$, so that $\dim_{\mathcal{H}}\Gamma = n$, even though $h$ need not be differentiable everywhere.
\end{exercise}

\begin{solution}
The map $g(x) = (x, h(x))$ satisfies $|g(x) - g(y)|^2 = |x - y|^2 + |h(x) - h(y)|^2 \leq (1 + L^2)|x - y|^2$, so it is $\sqrt{1 + L^2}$-Lipschitz, and its inverse, the projection
$\Gamma \to U$, is $1$-Lipschitz. By Proposition~\ref{geometric-measure-theory:proposition-hausdorff}(3) and Theorem~\ref{geometric-measure-theory:theorem-hausdorff-lebesgue},
$\lambda^n(U) = \mathcal{H}^n(U) \leq \mathcal{H}^n(\Gamma) \leq (1 + L^2)^{n/2}\mathcal{H}^n(U)$. Both bounds are positive and finite, so $\mathcal{H}^s(\Gamma)$ is $\infty$ for $s < n$ and $0$ for
$s > n$ by Proposition~\ref{geometric-measure-theory:proposition-hausdorff}(4). (By the general area formula of Remark~\ref{geometric-measure-theory:remark-area-coarea} and Rademacher's theorem,
$\mathcal{H}^n(\Gamma) = \int_U\sqrt{1 + |\nabla h|^2}$.)
\end{solution}

\begin{exercise}\label{geometric-measure-theory:exercise-torus-area}
For $R > r > 0$ let $x(u, v) = ((R + r\cos u)\cos v, (R + r\cos u)\sin v, r\sin u)$. Show that $x$ is injective with injective differential on $(0, 2\pi)^2$, that $Jx = r(R + r\cos u)$, and that
the torus $T = x(\RR^2)$ has $\mathcal{H}^2(T) = 4\pi^2Rr$.
\end{exercise}

\begin{solution}
The partial derivatives
\[\begin{aligned}
x_u &= (-r\sin u\cos v, -r\sin u\sin v, r\cos u), \\
x_v &= (-(R + r\cos u)\sin v, (R + r\cos u)\cos v, 0)
\end{aligned}\]
are orthogonal, of lengths $r$ and $R + r\cos u > 0$, so
$Jx = |x_u \times x_v| = r(R + r\cos u)$ and $Dx$ is injective. The point $x(u, v)$ determines $v$ modulo $2\pi$ (its projection to the plane $z = 0$ has argument $v$) and then $u$ modulo
$2\pi$, so $x$ is injective on $(0, 2\pi)^2$. Its image is $T$ minus the two circles $u = 0$ and $v = 0$, which have finite $\mathcal{H}^1$-measure and are therefore $\mathcal{H}^2$-null. By the area
formula, $\mathcal{H}^2(T) = \int_0^{2\pi}\int_0^{2\pi}r(R + r\cos u)\,du\,dv = 4\pi^2Rr$.
\end{solution}

\begin{exercise}\label{geometric-measure-theory:exercise-box-cantor}
Show directly from Example~\ref{geometric-measure-theory:example-cantor-dimension}(2), without Theorem~\ref{geometric-measure-theory:theorem-moran}, that the Cantor set has box-counting dimension $\log 2/\log 3$.
\end{exercise}

\begin{solution}
Let $3^{-k-1} \leq \delta < 3^{-k}$. The $2^{k+1}$ intervals of generation $k + 1$ have length $3^{-k-1} \leq \delta$ and cover $C$, so $N_\delta(C) \leq 2^{k+1}$. A set of diameter at most $\delta < 3^{-k}$
meets at most one of the $2^k$ intervals of generation $k$, and each of them contains points of $C$, so $N_\delta(C) \geq 2^k$. Since $k\log 3 < -\log\delta \leq (k + 1)\log 3$,
\[
\frac{k\log 2}{(k + 1)\log 3} \leq \frac{\log N_\delta(C)}{-\log\delta} \leq \frac{(k + 1)\log 2}{k\log 3} \qquad (k \geq 1),
\]
and both bounds tend to $\log 2/\log 3$ as $\delta \to 0$, that is, as $k \to \infty$.
\end{solution}

\begin{exercise}\label{geometric-measure-theory:exercise-cantor-sum}
Let $C$ be the Cantor set. Show that $C + C = \{x + y : x, y \in C\}$ is the interval $[0, 2]$. (Hint: $C \times C$ is the attractor of
Example~\ref{geometric-measure-theory:example-similarity-dimension}(3); project it as in Example~\ref{geometric-measure-theory:example-four-corner-projections}.) So the set $C \times C$, of dimension $\log 4/\log 3 < 2$, has a
projection of positive length, while its projections to the axes are copies of $C$, of length $0$.
\end{exercise}

\begin{solution}
Let $\pi(x, y) = x + y$ and $f_{a, b}(x, y) = (x/3 + a, y/3 + b)$ for $a, b \in \{0, \frac23\}$. Then $\pi \circ f_{a, b} = g_{a + b} \circ \pi$ with $g_c(t) = t/3 + c$, and $a + b$ takes the values
$0, \frac23, \frac43$. Since $C \times C = \bigcup_{a, b}f_{a, b}(C \times C)$, its image $\pi(C \times C) = C + C$ is a nonempty compact set with $C + C = g_0(C + C) \cup g_{2/3}(C + C) \cup g_{4/3}(C + C)$,
that is, the attractor of $\{g_0, g_{2/3}, g_{4/3}\}$. The interval $[0, 2]$ is also invariant: $g_0([0, 2]) \cup g_{2/3}([0, 2]) \cup g_{4/3}([0, 2]) = [0, \frac23] \cup [\frac23, \frac43] \cup [\frac43, 2]$. By
the uniqueness in Theorem~\ref{geometric-measure-theory:theorem-hutchinson}, $C + C = [0, 2]$. (More generally, by Marstrand's projection theorem a Borel set $A \subseteq \RR^2$ with $\dim_{\mathcal{H}}A > 1$ has
projections of positive length in almost every direction \cite{mattila-1995}.)
\end{solution}

\begin{exercise}\label{geometric-measure-theory:exercise-prescribed-dimension}
Show that for every $s \in [0, 1]$ there is a compact set $K \subseteq [0, 1]$ with $\dim_{\mathcal{H}}K = s$ and $0 < \mathcal{H}^s(K) < \infty$.
\end{exercise}

\begin{solution}
For $s = 0$ take a point, for $s = 1$ the interval $[0, 1]$ (Theorem~\ref{geometric-measure-theory:theorem-hausdorff-lebesgue}). For $0 < s < 1$ let $\rho = 2^{-1/s}$, so $0 < \rho < \frac12$, and consider
$f_1(x) = \rho x$, $f_2(x) = \rho x + 1 - \rho$. Since $f_1([0, 1]) \cup f_2([0, 1]) = [0, \rho] \cup [1 - \rho, 1] \subseteq [0, 1]$, the attractor $K$ lies in $[0, 1]$
(Theorem~\ref{geometric-measure-theory:theorem-hutchinson}(1)), and the pieces $f_1(K) \subseteq [0, \rho]$ and $f_2(K) \subseteq [1 - \rho, 1]$ are disjoint because $\rho < 1 - \rho$. The similarity dimension solves
$2\rho^t = 1$, so it is $s$, and Theorem~\ref{geometric-measure-theory:theorem-moran} applies.
\end{solution}

\begin{exercise}\label{geometric-measure-theory:exercise-natural-measure}
Let $K$ be the attractor of an iterated function system satisfying the strong separation condition, with similarity dimension $s$, and let $\mu$ be the measure constructed in the proof of
Theorem~\ref{geometric-measure-theory:theorem-moran}.
\begin{enumerate}
\item Show that $\mathcal{H}^s(K_w) = r_w^s\mathcal{H}^s(K)$ for every word $w$, and explain why this, together with $0 < \mathcal{H}^s(K) < \infty$, forces \eqref{geometric-measure-theory:equation-moran}.
\item Show that two Borel probability measures on $K$ which agree on all the sets $K_w$ are equal. Deduce that $\mu(B) = \mathcal{H}^s(B)/\mathcal{H}^s(K)$ for every Borel set $B \subseteq K$.
\end{enumerate}
\end{exercise}

\begin{solution}
(1) $f_w$ is a similarity with ratio $r_w$, so $\mathcal{H}^s(f_w(K)) = r_w^s\mathcal{H}^s(K)$. The sets $K_1, \ldots, K_N$ are disjoint compact sets with union $K$, so
$\mathcal{H}^s(K) = \sum_i\mathcal{H}^s(K_i) = \bigl(\sum_ir_i^s\bigr)\mathcal{H}^s(K)$, and dividing by $\mathcal{H}^s(K) \in (0, \infty)$ gives $\sum_ir_i^s = 1$. This is the heuristic of the introduction
to Section~\ref{geometric-measure-theory:section-self-similar}.

(2) Let $\mathcal{P}$ consist of $\emptyset$ and the sets $K_w$, including $K = K_\varnothing$. Two of the $K_w$ are either nested (if one word begins with the other, by
Theorem~\ref{geometric-measure-theory:theorem-hutchinson}(2)) or disjoint (if the words first differ at some position, by Step 1 of the proof of Theorem~\ref{geometric-measure-theory:theorem-moran}), so $\mathcal{P}$ is
closed under finite intersections. Every relatively open $G \subseteq K$ is a countable union of sets $K_w \subseteq G$: if $x \in G$ has address $\omega$, then $x \in K_{\omega|k}$ and
$\operatorname{diam}K_{\omega|k} \leq r^k\operatorname{diam}K$, so $K_{\omega|k} \subseteq G$ for $k$ large; and there are only countably many words. So $\sigma(\mathcal{P})$ contains the open subsets of
$K$, hence all Borel subsets, and Lemma~\ref{measure-theory:lemma-dynkin} shows that two probability measures agreeing on $\mathcal{P}$ agree on all Borel sets. Finally,
$B \mapsto \mathcal{H}^s(B)/\mathcal{H}^s(K)$ is a Borel probability measure on $K$ (Proposition~\ref{geometric-measure-theory:proposition-hausdorff}(1) and Theorem~\ref{geometric-measure-theory:theorem-moran}) which by (1)
gives the value $r_w^s = \mu(K_w)$ to every $K_w$.
\end{solution}

\begin{exercise}\label{geometric-measure-theory:exercise-sierpinski}
Let $S$ be the Sierpi\'nski gasket of a triangle $T$ with perimeter $P$ (Example~\ref{geometric-measure-theory:example-ifs}(3)). Show:
\begin{enumerate}
\item $\lambda^2(S) = 0$;
\item two distinct triangles $f_v(T)$, $f_w(T)$ with $|v| = |w|$ meet in at most one point;
\item $S$ contains the boundary $\partial T$, and more generally the boundaries of all triangles $f_w(T)$;
\item $\mathcal{H}^1(S) = \infty$, so that $1 \leq \dim_{\mathcal{H}}S \leq \log 3/\log 2$.
\end{enumerate}
\end{exercise}

\begin{solution}
(1) $F^k(T)$ is the union of the $3^k$ triangles $f_w(T)$, $|w| = k$, which are images of $T$ under similarities of ratio $2^{-k}$, so of area $4^{-k}\lambda^2(T)$ (Proposition~\ref{geometric-measure-theory:proposition-hausdorff}(3)
and Theorem~\ref{geometric-measure-theory:theorem-hausdorff-lebesgue}). Hence $\lambda^2(S) \leq \lambda^2(F^k(T)) \leq (\frac34)^k\lambda^2(T) \to 0$.

(2) For $i \neq j$ the corner triangles $f_i(T)$ and $f_j(T)$ meet only in the midpoint of the side $p_ip_j$. If $v \neq w$ first differ at position $j$ and $u$ is their common beginning of length
$j - 1$, then $f_v(T) \subseteq f_u(f_{v_j}(T))$ and $f_w(T) \subseteq f_u(f_{w_j}(T))$, and these meet in the image under the injective map $f_u$ of a single point.

(3) The side $[p_i, p_j]$ is the union of $[p_i, m]$ and $[m, p_j]$, $m$ its midpoint, which are sides of $f_i(T)$ and $f_j(T)$. So $\partial T \subseteq F(\partial T)$, and inductively
$\partial T \subseteq F(\partial T) \subseteq F^2(\partial T) \subseteq \cdots$, since $F$ preserves inclusions. As $F^k(\partial T) \subseteq F^k(T)$, $\partial T \subseteq \bigcap_kF^k(T) = S$
(Theorem~\ref{geometric-measure-theory:theorem-hutchinson}(1)). Then $f_w(\partial T) \subseteq f_w(S) \subseteq S$.

(4) A segment of length $\ell$ is isometric to an interval, so it has $\mathcal{H}^1$-measure $\ell$ (Theorem~\ref{geometric-measure-theory:theorem-hausdorff-lebesgue}). By (3), $S$ contains the boundaries of the $3^k$
triangles of generation $k$, each the union of three segments of total length $2^{-k}P$. By (2) these $3^{k+1}$ segments meet each other only in finitely many points, which are $\mathcal{H}^1$-null; so,
$\mathcal{H}^1$ being a measure on Borel sets, $\mathcal{H}^1(S) \geq 3^k2^{-k}P = (\frac32)^kP$ for every $k$. Hence $\mathcal{H}^1(S) = \infty$ and $\dim_{\mathcal{H}}S \geq 1$; the upper bound is
Proposition~\ref{geometric-measure-theory:proposition-similarity-upper}, with $3 \cdot 2^{-s} = 1$. (By Remark~\ref{geometric-measure-theory:remark-open-set-condition}, $\dim_{\mathcal{H}}S = \log 3/\log 2$.)
\end{solution}

\begin{exercise}\label{geometric-measure-theory:exercise-holder-graph}
Let $h \colon [0, 1] \to \RR$ satisfy $|h(x) - h(y)| \leq c|x - y|^\alpha$ for some $c > 0$ and $0 < \alpha \leq 1$ (a \emph{H\"older continuous} function), and let
$G = \{(x, h(x)) : x \in [0, 1]\}$ be its graph. Show that $1 \leq \dim_{\mathcal{H}}G \leq \overline{\dim}_BG \leq 2 - \alpha$.
\end{exercise}

\begin{solution}
Projection to the first coordinate is $1$-Lipschitz and maps $G$ onto $[0, 1]$, so $\mathcal{H}^1(G) \geq 1$ and $\dim_{\mathcal{H}}G \geq 1$. For $m \geq 1$ divide $[0, 1]$ into $m$ intervals $I_j$ of
length $\frac1m$. On $I_j$ the values of $h$ lie in an interval of length at most $cm^{-\alpha}$, which is covered by at most $cm^{1-\alpha} + 1$ intervals of length $\frac1m$; so the part of $G$ over
$I_j$ is covered by at most $cm^{1-\alpha} + 1$ squares of side $\frac1m$, and $G$ by at most $m(cm^{1-\alpha} + 1) \leq (c + 1)m^{2-\alpha}$ of them. The squares have diameter $\sqrt2/m$. Given
$0 < \delta < \sqrt2$, choose $m$ with $\sqrt2/(m + 1) \leq \delta < \sqrt2/m$; then $N_\delta(G) \leq (c + 1)(m + 1)^{2-\alpha}$ and $-\log\delta > \log(m/\sqrt2)$, so
\[
\frac{\log N_\delta(G)}{-\log\delta} \leq \frac{\log(c + 1) + (2 - \alpha)\log(m + 1)}{\log m - \log\sqrt2} \to 2 - \alpha \qquad (\delta \to 0).
\]
The middle inequality is Proposition~\ref{geometric-measure-theory:proposition-dimension-properties}(3). For $\alpha = 1$, a Lipschitz function, the graph has dimension $1$, in accordance with
Exercise~\ref{geometric-measure-theory:exercise-lipschitz-graph}; rough functions can have graphs of dimension close to $2$.
\end{solution}

\begin{exercise}\label{geometric-measure-theory:exercise-ball-derivative}
Let $g \colon \RR^n \to \RR$ be continuous, $n \geq 2$, and $\Phi(R) = \int_{B(0, R)}g\,d\lambda^n$. Show that $\Phi$ is differentiable on $(0, \infty)$ with
$\Phi^{\prime}(R) = \int_{|x| = R}g\,d\mathcal{H}^{n-1}$, and that the averages $\frac1{n\omega_nR^{n-1}}\int_{|x| = R}g\,d\mathcal{H}^{n-1}$ tend to $g(0)$ as $R \to 0$.
\end{exercise}

\begin{solution}
Apply Corollary~\ref{geometric-measure-theory:corollary-polar} to the positive and negative parts of $g\mathbf{1}_{B(0, R)}$, which are bounded Borel functions vanishing outside a bounded set; subtracting,
$\Phi(R) = \int_0^R\eta(r)\,dr$ with $\eta(r) = r^{n-1}\int_{S^{n-1}}g(r\theta)\,d\mathcal{H}^{n-1}(\theta)$. As $r \to r_0$, $g(r\theta) \to g(r_0\theta)$ uniformly in $\theta \in S^{n-1}$, by uniform
continuity of $g$ on a compact ball, and $\mathcal{H}^{n-1}(S^{n-1}) = n\omega_n < \infty$; so $\eta$ is continuous on $[0, \infty)$, and $\Phi^{\prime}(R) = \eta(R)$ by the fundamental theorem of calculus
(Theorem~\ref{real-analysis:theorem-ftc}). By the scaling in the proof of Corollary~\ref{geometric-measure-theory:corollary-polar}, $\eta(R) = \int_{|x| = R}g\,d\mathcal{H}^{n-1}$. Finally
$\eta(R)/(n\omega_nR^{n-1}) - g(0) = \frac{1}{n\omega_n}\int_{S^{n-1}}(g(R\theta) - g(0))\,d\mathcal{H}^{n-1}(\theta)$, which is at most $\sup_{|x| = R}|g(x) - g(0)| \to 0$ in absolute value.
\end{solution}

\begin{exercise}\label{geometric-measure-theory:exercise-pappus-cone}
Use the Pappus--Guldin theorems to compute the volume and the lateral area of the cone of revolution with base radius $R$ and height $h$, obtained by rotating the triangle with vertices $(0, 0)$, $(R, 0)$,
$(0, h)$ of the $(r, z)$-plane around the $z$-axis. (The centroid of a triangle is the average of its vertices.)
\end{exercise}

\begin{solution}
The triangle has area $\frac12Rh$ and centroid at distance $\bar r = \frac13(0 + R + 0) = \frac R3$ from the axis, so the volume is $2\pi\frac R3 \cdot \frac12Rh = \frac13\pi R^2h$, by
Corollary~\ref{measure-theory:corollary-pappus-volume}. The lateral surface comes from the segment from $(R, 0)$ to $(0, h)$, of length $\sqrt{R^2 + h^2}$, whose centroid is its midpoint, at distance $\frac R2$; by
Proposition~\ref{geometric-measure-theory:proposition-pappus-area} the area is $2\pi\frac R2\sqrt{R^2 + h^2} = \pi R\sqrt{R^2 + h^2}$. (That the centroid of a triangle is the average of its vertices follows by applying an affine map
taking it to the triangle $(0, 0)$, $(1, 0)$, $(0, 1)$, which changes all areas by the same factor, and computing $\int_0^1\int_0^{1-r}r\,dz\,dr = \frac16 = \frac13 \cdot \frac12$.)
\end{solution}
